\documentclass[11pt,a4paper]{article}

%===============================================================================
% PACKAGES
%===============================================================================
\usepackage[a4paper,margin=2.5cm]{geometry}

\usepackage{amsmath,amssymb,amsfonts,mathtools}
\usepackage{bm}
\usepackage{bbm}
\usepackage{graphicx}
\usepackage{booktabs}
\usepackage{enumitem}
\usepackage{xcolor}
\usepackage{caption}
\usepackage{subcaption}
\usepackage{hyperref}
\usepackage{algorithm}
\usepackage{algpseudocode}
\usepackage{float}
\hypersetup{
    colorlinks=true,
    linkcolor=blue!50!black,
    citecolor=blue!50!black,
    urlcolor=blue!50!black
}


%===============================================================================
% MACROS
%===============================================================================

%------------------------------------------------------------------------------
% General sets and dimensions
%------------------------------------------------------------------------------
\newcommand{\R}{\mathbb{R}}
\newcommand{\Nip}{n_{\mathrm{ip}}}       % full number of integration points
\newcommand{\Nman}{N_{\mathrm{M}}}       % number of sampled manifold states
\newcommand{\Ncand}{n_{\mathrm{c}}}      % number of current candidate points
\newcommand{\Nred}{n_{\mathrm{red}}}     % number of final retained points
\newcommand{\Nloc}{n_{\mathrm{loc}}}     % number of local constraints
\newcommand{\Nlat}{r_{\mathrm{M}}}       % latent dimension

%------------------------------------------------------------------------------
% Latent coordinates and sampled manifold
%------------------------------------------------------------------------------
\newcommand{\qM}{\bm q_{\mathrm{M}}}
\newcommand{\qMj}{\bm q_{\mathrm{M}}^{(j)}}
\newcommand{\Qman}{\bm Q_{\mathrm{M}}}
\newcommand{\Mtrain}{\mathcal{M}_{\mathrm{tr}}}

%------------------------------------------------------------------------------
% Integration points and weights
%------------------------------------------------------------------------------
\newcommand{\Zfull}{\mathcal{Z}_{\mathrm{FE}}}
\newcommand{\Zini}{\mathcal{Z}_{0}}
\newcommand{\Zred}{\mathcal{Z}_{\mathrm{red}}}
\newcommand{\zini}{\bm z_{0}}
\newcommand{\zred}{\bm z_{\mathrm{red}}}

\newcommand{\wFE}{\bm w^{\mathrm{FE}}}
\newcommand{\wini}{\bm w_{0}}
\newcommand{\wred}{\bm w_{\mathrm{red}}}
\newcommand{\wadapt}{\bm w}
\newcommand{\Wadapt}{\bm W}
\newcommand{\Wold}{\bm W^{\mathrm{old}}}
\newcommand{\Wnew}{\bm W^{\mathrm{new}}}
\newcommand{\DeltaW}{\Delta \bm W}

%------------------------------------------------------------------------------
% Local bases, constraints, and integrand matrices
%------------------------------------------------------------------------------
\newcommand{\Aint}{\bm A}
\newcommand{\Aj}{\bm A_j}
\newcommand{\Uj}{\bm U_j}
\newcommand{\bj}{\bm b_j}
\newcommand{\Ufixed}{\bm U_{\mathrm{fix}}}
\newcommand{\Qproj}{\bm Q}
\newcommand{\bone}{\bm 1}

%------------------------------------------------------------------------------
% Graph quantities
%------------------------------------------------------------------------------
\newcommand{\Kgraph}{\bm K_{\mathrm{G}}}
\newcommand{\Ggraph}{\mathcal{G}_{\mathrm{M}}}
\newcommand{\Egraph}{\mathcal{E}_{\mathrm{M}}}
\newcommand{\alphaG}{\alpha_{\mathrm{G}}}

%------------------------------------------------------------------------------
% Pruning and active-set notation
%------------------------------------------------------------------------------
\newcommand{\Iact}{\mathcal{I}_{\mathrm{act}}}
\newcommand{\Ielim}{\mathcal{I}_{\mathrm{elim}}}
\newcommand{\Ipos}{\mathcal{I}_{+}}
\newcommand{\Dact}{\mathcal{D}}
\newcommand{\Dj}{\mathcal{D}_j}

%------------------------------------------------------------------------------
% Operators
%------------------------------------------------------------------------------
\DeclareMathOperator{\vecop}{vec}
\DeclareMathOperator{\range}{Range}
\DeclareMathOperator{\rank}{rank}
\DeclareMathOperator{\trace}{tr}
\DeclareMathOperator{\diag}{diag}
\DeclareMathOperator{\supp}{supp}
\newcommand{\Qall}{\bm Q_{\mathrm{all}}}



%===============================================================================
% ADDITIONAL MACROS FOR LOCAL EXACTNESS SYSTEMS SECTION
%===============================================================================

%------------------------------------------------------------------------------
% Full and projected local bases
%------------------------------------------------------------------------------
\newcommand{\Ufull}{\bm U_{\mathrm{full}}}
\newcommand{\Ufullj}{\bm U_{\mathrm{full},j}}

%------------------------------------------------------------------------------
% FE-weighted scaling matrix
%------------------------------------------------------------------------------
\newcommand{\Sfe}{\bm S_{\mathrm{FE}}}

%------------------------------------------------------------------------------
% Projected / filtered integrand matrices
%------------------------------------------------------------------------------
\newcommand{\Aproj}{\bm A^{\mathrm{proj}}}
\newcommand{\Ahat}{\widehat{\bm A}}

%------------------------------------------------------------------------------
% Aggregate candidate weights
%------------------------------------------------------------------------------
\newcommand{\Ri}{R_i}

%------------------------------------------------------------------------------
% Graph-based latent quantities
%------------------------------------------------------------------------------
\newcommand{\Qone}{Q_1}
\newcommand{\Qtwo}{Q_2}

%------------------------------------------------------------------------------
% Manifold samples
%------------------------------------------------------------------------------
\newcommand{\qMi}{\bm q_{\mathrm{M}}^{(i)}}

%------------------------------------------------------------------------------
% Candidate support sets
%------------------------------------------------------------------------------
\newcommand{\Zcand}{\mathcal{Z}_{\mathrm{cand}}}

%------------------------------------------------------------------------------
% Optional fixed elastic scaffold
%------------------------------------------------------------------------------
\newcommand{\Uel}{\bm U_{\mathrm{el}}}

%------------------------------------------------------------------------------
% Sample-dependent adaptive weights
%------------------------------------------------------------------------------
\newcommand{\wadaptj}{\bm w^{(j)}}

%------------------------------------------------------------------------------
% Sample-dependent integrand matrices
%------------------------------------------------------------------------------
\newcommand{\Aintj}{\bm A_j}


%===============================================================================
% ADDITIONAL MACROS FOR GRAPH-LAPLACIAN SECTION
%===============================================================================

\newcommand{\Tgraph}{\mathcal{T}_{\mathrm{G}}}
\newcommand{\xG}{\bm x^{\mathrm{G}}}
\newcommand{\varphivec}{\bm\varphi}
\newcommand{\SG}{\mathcal{S}_{\mathrm{G}}}
\newcommand{\zvec}{\bm z}
\newcommand{\zold}{\bm z^{\mathrm{old}}}
\newcommand{\Hgraph}{\bm H_{\mathrm{G}}}
\newcommand{\Iid}{\bm I}
\newcommand{\lambdaB}{\lambda_{\mathrm{B}}}
\newcommand{\ellB}{\ell_{\mathrm{B}}}
\newcommand{\rhoB}{\rho}



%===============================================================================
% ADDITIONAL MACROS FOR PRUNING SECTION
%===============================================================================

\newcommand{\wnew}{\bm w^{\mathrm{new}}}
\newcommand{\wold}{\bm w^{\mathrm{old}}}
\newcommand{\Wred}{\bm W_{\mathrm{red}}}
\newcommand{\rmax}{r_{\max}}
\newcommand{\nstop}{n_{\mathrm{stop}}}
\newcommand{\ntarget}{n_{\mathrm{target}}}


%===============================================================================
% ADDITIONAL MACROS FOR NO-ENFORCEMENT SECTION
%===============================================================================

\newcommand{\Irem}{\mathcal{I}_{\mathrm{rem}}}
\newcommand{\wnewj}{\bm w_{\mathrm{new}}^{(j)}}
\newcommand{\woldj}{\bm w_{\mathrm{old}}^{(j)}}
\newcommand{\lambdaj}{\bm\lambda_j}
\newcommand{\Lloc}{\mathcal{L}}


%===============================================================================
% ADDITIONAL MACROS FOR LOCAL ACTIVE-SET SECTION
%===============================================================================

\newcommand{\Ifree}{\mathcal{I}_{\mathrm{free}}}
\newcommand{\Ifreej}{\mathcal{I}_{\mathrm{free},j}}
\newcommand{\Nneg}{\mathcal{N}}
\newcommand{\Nnegj}{\mathcal{N}_j}
\newcommand{\Afeasible}{\mathcal{A}}
\newcommand{\Ffeasible}{\mathcal{F}}


%===============================================================================
% ADDITIONAL MACROS FOR GRAPH-REGULARIZED STAGE
%===============================================================================

\newcommand{\zp}{\bm z_{\mathrm{p}}}
\newcommand{\Nmat}{\bm N}
\newcommand{\yvec}{\bm y}
\newcommand{\Hred}{\bm H_{\mathrm{red}}}
\newcommand{\fred}{\bm f_{\mathrm{red}}}
\newcommand{\gvec}{\bm g}
\newcommand{\epsneg}{\varepsilon_{\mathrm{neg}}}
\newcommand{\kmax}{k_{\max}}


%===============================================================================
% ADDITIONAL MACROS FOR WEIGHT REGRESSION SECTION
%===============================================================================

\newcommand{\PhiRBF}{\bm\Phi}
\newcommand{\gammai}{\bm\gamma_i}
\newcommand{\etai}{\bm\eta_i}
\newcommand{\ci}{\bm c_i}
\newcommand{\Di}{\bm D_i}
\newcommand{\wfield}{\bm w_i}



%------------------------------------------------------------------------------
% MATLAB/file references
%------------------------------------------------------------------------------
\newcommand{\MATLAB}[1]{%
\begin{quote}
\small\ttfamily #1
\end{quote}
}

%===============================================================================
% DOCUMENT
%===============================================================================

\title{
Bivariate Manifold-Adaptive Weights Empirical Cubature Method\\
\large Construction, Graph-Regularized Pruning, and Weight Regression
}

\author{Joaqu\'in A. Hern\'andez Ortega}

\date{May 2026}

\begin{document}

\maketitle

\begin{abstract}

See shorter version in \href{/home/joaquin/Desktop/CURRENT_TASKS/MATLAB_CODES/TESTING_PROBLEMS_FEHROM/115_PAPER_MAW_ECM/LATEX/02_MAW_ECM.pdf}{02\_MAW\_ECM.pdf}.


This document develops, from first principles, the formulation implemented in
\texttt{MAW\_ECM\_genV1.m}. The method is a manifold-adaptive extension of the
Empirical Cubature Method (ECM), constructed for nonlinear manifold reduced-order
models in which the relevant reduced states lie on a low-dimensional latent
manifold.

The formulation originates from the idea of combining a common reduced
integration support with manifold-dependent cubature weights. The resulting
method may be interpreted as a manifold counterpart of subspace-adaptive
cubature constructions such as SAW--ECM, in which the support is shared across a
family of reduced representations while the weights are allowed to vary.
Accordingly, the central object in MAW--ECM is not a single cubature vector, but
a family of adaptive weights
\[
        \wadapt=\wadapt(\qM),
\]
where $\qM$ denotes the latent coordinates of the nonlinear manifold.

At the conceptual level, the method may be viewed as a manifold-adaptive
sparsification problem: one seeks the smallest common support capable of
reproducing, at every sampled manifold state, a collection of local exactness
constraints associated with reduced residuals and outputs. However, the direct
solution of this global support-selection problem is computationally intractable
because it combines combinatorial sparsity, positivity, local exactness, and
manifold regularity constraints in a single coupled formulation.

The implemented MAW--ECM construction therefore does not solve the original
sparsification problem directly. Instead, it proceeds through a feasible pruning
strategy starting from an already admissible fixed-weight ECM rule. The initial
ECM support is progressively reduced by removing candidate integration points
one at a time and redistributing the remaining weights over the manifold while
preserving the local exactness conditions. In this sense, sparsification defines
the ideal mathematical objective, whereas pruning defines the practical
algorithmic route.

The redistribution process operates in two regimes. First, an inexpensive
least-change stage attempts to eliminate points without explicit positivity
enforcement. Once this becomes impossible, the method switches to a
graph-regularized active-set redistribution stage in which positivity is
enforced while neighboring manifold states are coupled through a graph
Laplacian. This graph coupling is introduced to moderate the loss of regularity
caused by local active-set transitions in the adaptive weights.

The retained integration points are common to all sampled manifold states,
whereas the associated weights vary over the latent manifold. Consequently, the
method lies between a fully global ECM rule, where both support and weights are
fixed, and a fully local cubature strategy, where both support and weights would
be recomputed independently at each state.

The present implementation is specialized to a bivariate latent manifold, in
which the sampled states lie on a structured two-dimensional grid. This
assumption is used primarily in the construction of the graph regularization
operator and in the visualization of the adaptive weight fields. The algebraic
structure of MAW--ECM itself is more general and extends naturally to latent
manifolds of dimension $\Nlat>2$, provided that a suitable neighborhood graph or
discrete manifold connectivity is available.
\end{abstract}

\tableofcontents

%===============================================================================
\section{Purpose and scope of the document}
%===============================================================================

The goal of this note is to fix the mathematical and algorithmic structure of
the manifold-adaptive weights ECM, abbreviated here as MAW--ECM. The document is
not intended to be inserted directly into a paper. Its purpose is instead to
make explicit the ideas implemented in the MATLAB routine
\[
        \texttt{MAW\_ECM\_genV1.m},
\]
and in the auxiliary routines that it invokes.

The central question addressed by the method is the following. Suppose that a
fixed ECM rule has already selected a set of integration points
\[
        \Zini \subset \Zfull := \{1,\ldots,\Nip\},
\]
with associated positive weights $\wini$. This rule is globally valid over the
training manifold, but it may still contain more points than appear necessary
from the viewpoint of the latent dimension. Can the support be reduced further
if the weights are allowed to vary with the manifold coordinate?

MAW--ECM answers this question by replacing a single fixed cubature vector by a
matrix of adaptive weights,
\[
        \Wadapt =
        \begin{bmatrix}
        | & | &        & | \\
        \wadapt^{(1)} & \wadapt^{(2)} & \cdots & \wadapt^{(\Nman)} \\
        | & | &        & |
        \end{bmatrix}
        \in \R^{\Ncand \times \Nman}.
\]
Here $\Ncand = |\Zini|$ is the number of current candidate points and $\Nman$ is
the number of sampled manifold states. The $j$-th column
$\wadapt^{(j)}$ contains the cubature weights associated with the latent sample
$\qMj$.

The support is pruned globally: removing one integration point means removing
one row of $\Wadapt$ for all sampled states. The weights, however, are adapted
locally in the manifold:
\[
        \wadapt^{(j)} \equiv \wadapt(\qMj).
\]
The final online approximation is obtained by fitting a regression model to the
columns of $\Wadapt$, so that the weights can be evaluated at a new latent
coordinate $\qM$.

%===============================================================================
\section{Relation with classical ECM}
%===============================================================================

Let $\Zfull=\{1,\ldots,\Nip\}$ denote the full set of finite element integration
points and let $\wFE\in\R^{\Nip}$ collect the corresponding finite element
quadrature weights. For a state-dependent scalar integrand
$f_i(\qM)$ evaluated at integration point $i$, the full quadrature approximation
reads
\[
        I(\qM)
        =
        \sum_{i=1}^{\Nip} f_i(\qM)\,w^{\mathrm{FE}}_i .
\]
For several scalar integrands collected in a matrix
\[
        \Aint(\qM)\in\R^{\Nip\times m},
\]
the full integral vector is
\[
        \bm b(\qM)
        =
        \Aint(\qM)^T \wFE .
\]

Classical ECM constructs a fixed subset $\Zred\subset\Zfull$ and a fixed vector
of positive weights $\wred$ such that
\[
        \Aint_{\Zred}({\qM}_j)^T\wred
        \approx
        \Aint({\qM}_j)^T\wFE,
        \qquad
        j=1,\ldots,\Nman .
\]
The same weights must be valid for all sampled states. Consequently, the fixed
rule must represent the union of all relevant integrand configurations over the
training manifold. This global requirement is precisely what limits the
compressibility of a fixed-weight cubature formula.

MAW--ECM modifies only one aspect of this construction. The support remains
fixed after the offline stage, but the weights are allowed to vary with the
latent state:
\[
        \Aint_{\Zred}(\qM)^T\wred(\qM)
        \approx
        \Aint(\qM)^T\wFE .
\]
Thus, the method does not abandon the ECM philosophy. It keeps the idea of a
sparse positive cubature rule defined on selected finite element integration
points. The difference is that the cubature weights become functions on the
reduced manifold.


%===============================================================================
\section{Introduction: from adaptive cubature supports to manifold-adaptive weights}
\label{sec:introduction_history}
%===============================================================================

The Manifold-Adaptive Weights Empirical Cubature Method (MAW--ECM) was not
conceived as a direct regression strategy for cubature weights, nor as a purely
heuristic modification of a standard ECM rule. Its formulation emerged from the
attempt to reconcile three requirements that are individually natural but
difficult to satisfy simultaneously:
\begin{enumerate}[label=(\roman*),leftmargin=1.1cm]
    \item the cubature rule should use a common reduced set of integration
    points over a family of reduced states;
    \item the cubature weights should be allowed to adapt to the position on a
    nonlinear manifold;
    \item the resulting weight fields should be sufficiently regular to be used
    in an online hyper-reduced model.
\end{enumerate}
The first requirement is inherited from subspace-adaptive cubature ideas, the
second from the manifold character of the reduced model, and the third from the
fact that the adaptive weights are not merely postprocessed quantities: they
enter the online evaluation of residuals, outputs, and, in nonlinear settings,
their consistent linearizations.

%-------------------------------------------------------------------------------
\subsection{Roots in SAW--ECM: common support with adaptive weights}
%-------------------------------------------------------------------------------

The first conceptual root of MAW--ECM is the Subspace Adaptive Weights Empirical
Cubature Method (SAW--ECM). In SAW--ECM, one considers several reduced
subspaces, indexed by $s=1,\ldots,n_s$, and associates with each of them a
cubature weight vector
\[
        \bm\zeta^{(s)}\in\R^{\Nip}.
\]
The key point is that sparsity is not imposed independently on each
$\bm\zeta^{(s)}$. Instead, sparsity is imposed on the union of their supports.
At the formal level, this can be expressed through an objective involving the
zero-norm of the aggregate weight vector,
\[
        \left\|
        \sum_{s=1}^{n_s}\bm\zeta^{(s)}
        \right\|_0.
\]
Thus, all subspaces share a common set of selected integration points, while
the weights are allowed to depend on the subspace.

This idea is central for MAW--ECM. The method keeps the same distinction:
\[
        \text{common support}
        \qquad\text{versus}\qquad
        \text{state-dependent weights}.
\]
What changes is the indexing set. SAW--ECM adapts the weights over a finite
collection of reduced subspaces. MAW--ECM adapts the weights over a sampled
nonlinear manifold of latent states.

Let
\[
        \Mtrain=\{\qMj\}_{j=1}^{\Nman}
\]
be the sampled training manifold. At each sampled state $\qMj$, one introduces
a cubature weight vector
\[
        \bm\zeta^{(j)}\in\R^{\Nip}.
\]
The direct manifold analogue of the SAW--ECM support-union objective is then
\[
        \left\|
        \sum_{j=1}^{\Nman}\bm\zeta^{(j)}
        \right\|_0.
        \label{eq:intro_support_union_full}
\]
The support of this aggregate vector identifies the integration points that are
used somewhere on the manifold. Since the weights are nonnegative, the support
of the sum is precisely the union of the supports of the individual rules.

%-------------------------------------------------------------------------------
\subsection{The full manifold-adaptive sparsification problem}
%-------------------------------------------------------------------------------

For each sampled manifold state $\qMj$, let
\[
        \Uj\in\R^{\Nip\times r_j},
        \qquad
        \bj\in\R^{r_j},
\]
denote the local exactness operator and the corresponding vector of exact
integrals. The constraints
\[
        {\Uj}^T\bm\zeta^{(j)}=\bj,
        \qquad
        j=1,\ldots,\Nman,
        \label{eq:intro_full_exactness}
\]
represent the quantities that the cubature rule must reproduce at each
manifold point. In a manifold HROM, these quantities may include tangent-space
projected residual contributions, output quantities such as reactions or
homogenized stresses, and the volume constraint. The latter prevents the
degenerate zero-weight solution and preserves the total measure represented by
the quadrature rule.

A first formal statement of the manifold-adaptive sparsification problem is
therefore
\[
\begin{aligned}
\min_{\bm\zeta^{(1)},\ldots,\bm\zeta^{(\Nman)}}
&\quad
\left\|
\sum_{j=1}^{\Nman}\bm\zeta^{(j)}
\right\|_0
\\
\text{subject to}
&\quad
{\Uj}^T\bm\zeta^{(j)}=\bj,
\qquad
j=1,\ldots,\Nman,
\\
&\quad
\bm\zeta^{(j)}\geq 0,
\qquad
j=1,\ldots,\Nman.
\end{aligned}
\label{eq:intro_full_maw_l0}
\]
This is the original, idealized MAW--ECM problem. It is the exact analogue of
the SAW--ECM sparsification principle, but with the discrete subspace index
replaced by a sampled manifold index.

However, this formulation is still incomplete for a manifold method. Unlike a
finite collection of unrelated subspaces, the samples $\qMj$ lie on a
low-dimensional latent domain. The weights are ultimately intended to define a
function
\[
        \qM\mapsto\wadapt(\qM),
\]
to be evaluated online by interpolation or regression. Therefore, the vectors
$\bm\zeta^{(j)}$ should not vary arbitrarily from one sampled state to its
neighbors.

This leads to the smoothness-augmented formal problem
\[
\begin{aligned}
\min_{\bm\zeta^{(1)},\ldots,\bm\zeta^{(\Nman)}}
&\quad
\left\|
\sum_{j=1}^{\Nman}\bm\zeta^{(j)}
\right\|_0
+
\frac{\alphaG}{2}
\sum_{i=1}^{\Nip}
{\bm\zeta_i}^T\Kgraph\bm\zeta_i
\\
\text{subject to}
&\quad
{\Uj}^T\bm\zeta^{(j)}=\bj,
\qquad
j=1,\ldots,\Nman,
\\
&\quad
\bm\zeta^{(j)}\geq 0,
\qquad
j=1,\ldots,\Nman.
\end{aligned}
\label{eq:intro_full_smooth_l0}
\]
Here $\bm\zeta_i\in\R^{\Nman}$ collects the values, over the sampled manifold,
of the weight associated with integration point $i$, and $\Kgraph$ is a graph
Laplacian on the sampled latent states. The graph term penalizes lack of
regularity of the weight fields over the manifold.

Problem~\eqref{eq:intro_full_smooth_l0} is the full conceptual formulation of
MAW--ECM. It contains, in a single statement, the desired common support, local
exactness, positivity, and manifold regularity.

%-------------------------------------------------------------------------------
\subsection{Why the full formulation is not used directly}
%-------------------------------------------------------------------------------

Although \eqref{eq:intro_full_smooth_l0} is conceptually useful, it is not a
practical algorithm. Its difficulty is not merely that it contains many
unknowns. The fundamental obstacle is the combination of a combinatorial support
selection term with equality constraints, positivity constraints, and a global
manifold-smoothness penalty.

The zero-norm term encodes the discrete decision of which integration points
are retained. The equality constraints encode the exactness of the reduced
integrals at every sampled manifold state. The positivity constraints impose
admissibility of the cubature weights. The graph term couples the sampled
states and destroys the complete separability that one would otherwise exploit
node by node.

Thus, a direct solution of \eqref{eq:intro_full_smooth_l0} would require
solving, in one monolithic step, the support-selection problem and the
manifold-coupled constrained redistribution problem. This is not the route
taken here.

The purpose of \eqref{eq:intro_full_smooth_l0} is instead to clarify the target
structure of the method. It identifies what the algorithm should approximate:
a sparse common support, with nonnegative weights that vary over the manifold
while preserving the local cubature constraints.

%-------------------------------------------------------------------------------
\subsection{Restriction to an initial ECM support}
%-------------------------------------------------------------------------------

The first practical simplification is to avoid working on the full set of FE
integration points. MAW--ECM starts from an already feasible fixed ECM rule,
with support
\[
        \Zini\subset\Zfull,
        \qquad
        |\Zini|=\Ncand,
\]
and weights
\[
        \wini\in\R^{\Ncand}.
\]
The adaptive weights are then defined only on this candidate support:
\[
        \wadapt^{(j)}\in\R^{\Ncand},
        \qquad
        j=1,\ldots,\Nman.
\]
This restriction is not merely a computational convenience. It reflects the
constructive philosophy of the method: the initial ECM rule provides a feasible
global integration rule, and MAW--ECM attempts to prune it further by exploiting
state-dependent weights.

The initial adaptive rule is obtained by replication:
\[
        \wadapt^{(j)}=\wini,
        \qquad
        j=1,\ldots,\Nman.
        \label{eq:intro_initial_replication}
\]
Thus, before pruning begins, the adaptive rule coincides with the fixed ECM
rule at every sampled manifold state. The weight field is constant over the
manifold.

The question becomes whether rows of the matrix
\[
        \Wadapt =
        [\wadapt^{(1)},\ldots,\wadapt^{(\Nman)}]
        \in\R^{\Ncand\times\Nman}
\]
can be progressively eliminated while preserving feasibility of the local
systems.

%-------------------------------------------------------------------------------
\subsection{CECM influence: pruning instead of direct sparsification}
%-------------------------------------------------------------------------------

The second practical simplification is inherited from the Continuous Empirical
Cubature Method (CECM). The relevant lesson from CECM is not that the same
problem is being solved, but that a difficult sparse cubature construction can
be approached by a controlled sequence of pruning operations starting from a
feasible rule.

In the present context, one does not attempt to solve the zero-norm problem
\eqref{eq:intro_full_smooth_l0} directly. Instead, one proceeds iteratively.
At a given pruning stage, a candidate integration point is tentatively removed
from the current support. The corresponding row of $\Wadapt$ is forced to
vanish for all sampled states. The remaining weights are then redistributed so
that the local exactness constraints remain satisfied.

If such a redistribution is feasible, the elimination is accepted. If it is not
feasible, the candidate is rejected and another candidate is tested. The pruning
sequence stops when no further feasible removal is found, or when a prescribed
stopping threshold is reached.

This changes the nature of the problem. The original support-selection problem
is replaced by a sequence of feasibility-preserving local decisions:
\[
        \text{remove one point}
        \quad\longrightarrow\quad
        \text{restore feasible weights}
        \quad\longrightarrow\quad
        \text{accept or reject}.
\]
Sparsity is therefore increased monotonically, but only through admissible
steps.

%-------------------------------------------------------------------------------
\subsection{From the full global step to the implemented two-regime algorithm}
%-------------------------------------------------------------------------------

The formal smoothness-augmented single-step redistribution problem would seek,
after removal of a candidate point, a new weight matrix $\Wnew$ minimizing a
least-change term plus a graph-smoothness term, subject to local exactness and
positivity:
\[
\begin{aligned}
\min_{\Wnew}
&\quad
\frac{1}{2}\|\Wnew-\Wold\|_F^2
+
\frac{\alphaG}{2}
\trace\left[
(\Wnew-\Wold)\Kgraph(\Wnew-\Wold)^T
\right]
\\
\text{subject to}
&\quad
{\Uj}^T{\wnew}^{(j)}=\bj,
\qquad
j=1,\ldots,\Nman,
\\
&\quad
{\wnew}^{(j)}\geq 0,
\qquad
j=1,\ldots,\Nman,
\\
&\quad
W^{\mathrm{new}}_{pj}=0,
\qquad
j=1,\ldots,\Nman .
\end{aligned}
\label{eq:intro_single_step_full}
\]
This is the clean mathematical form of one graph-regularized pruning step.

Nevertheless, solving \eqref{eq:intro_single_step_full} at every pruning
iteration is unnecessary during the early part of the reduction. Numerical
experience shows that many points can be removed by a simpler equality-only
least-change redistribution. In that regime, positivity is not yet active, and
the local systems vary smoothly enough over the manifold that smoothness is
often obtained without explicitly imposing the graph penalty.

This observation shaped the implemented algorithm. The current MAW--ECM
construction uses two regimes:
\begin{enumerate}[label=(\roman*),leftmargin=1.1cm]
    \item a no-enforcement stage, in which candidate eliminations are attempted
    through inexpensive least-change redistributions without explicit
    positivity enforcement;
    \item a graph-regularized active-set stage, activated once positivity
    enforcement becomes necessary.
\end{enumerate}
The first regime is efficient and tends to preserve regularity because it avoids
active-set clamping. The second regime is more expensive but necessary near the
feasibility boundary, where nonnegative weights can no longer be recovered by
unconstrained redistribution alone.

In the current implementation, once explicit enforcement is required, the
graph-based version is the intended route. The purely local active-set variant
is useful historically and diagnostically, but the working MAW--ECM formulation
couples the manifold samples through the graph Laplacian whenever positivity
must be enforced.

%-------------------------------------------------------------------------------
\subsection{Role of positivity and active sets}
%-------------------------------------------------------------------------------

The appearance of negative weights is not a regression failure. It is a
structural event in the constrained pruning process. As the support shrinks, the
remaining weights must satisfy the same local exactness constraints with fewer
degrees of freedom. The equality-only least-change solution may then leave the
positive orthant.

When this happens, nonnegativity is enforced by an active-set strategy. A
negative weight is clamped to zero and the redistribution problem is solved
again on the remaining free variables. In MAW--ECM this clamping is local in
the manifold direction:
\[
        W_{ij}=0
        \quad\text{may be imposed at one manifold node }j,
\]
without imposing the same condition at all other nodes. This is distinct from
pruning, which is global:
\[
        W_{pj}=0
        \quad\text{for all }j=1,\ldots,\Nman.
\]
Keeping these two mechanisms separate is essential. Pruning removes an
integration point from the common support. Active-set clamping only restores
nonnegativity at specific entries of the weight field.

This distinction also explains why the resulting map
\[
        \qM\mapsto\wred(\qM)
\]
may be less regular than the underlying mechanical manifold. The weight map is
the solution map of a constrained optimization problem. It can change
nonsmoothly when active sets change. Graph regularization does not remove this
fact; it only selects, among feasible redistributions, those that introduce
less roughness across neighboring latent states.

%-------------------------------------------------------------------------------
\subsection{Current formulation}
%-------------------------------------------------------------------------------

The final form of the implemented MAW--ECM method can therefore be summarized as
follows.

First, a fixed ECM rule provides an admissible initial support and weights.
Second, this rule is lifted to the manifold by assigning the same initial
weights to every sampled latent state. Third, local exactness systems are built
at each sampled state, including the constraints required by the HROM residual
and by the selected outputs. Fourth, points are removed one by one. Each removal
is accepted only if the remaining weights can be redistributed so as to satisfy
the local constraints and, when required, positivity. Fifth, once the final
support has been obtained, the sampled adaptive weights are regressed as
functions of the latent coordinates.

Thus, the method realizes the original support-union idea of SAW--ECM and the
constructive pruning philosophy associated with CECM, but adapts both to the
specific needs of a nonlinear manifold HROM. The result is a fixed-support,
manifold-adaptive cubature rule:
\[
        \int_{\Omega_0} f(x,\qM)\,d\Omega
        \approx
        \sum_{i\in\Zred}
        w_i(\qM)\,
        f(x_i,\qM).
\]
The support $\Zred$ is selected offline and remains fixed online. The weights
are functions of the latent coordinates. The graph regularization is not an
independent modeling assumption, but a practical mechanism introduced when
positivity enforcement would otherwise generate irregular weight fields.

This historical route is important because it clarifies what MAW--ECM is and
what it is not. It is not a black-box regression of quadrature weights. It is
not a fully local cubature rule with state-dependent supports. It is not the
direct solution of the ideal zero-norm formulation. Rather, it is a feasible
pruning approximation to that formulation, built from a fixed ECM rule, enriched
with manifold-dependent weights, and regularized only as much as needed to make
the adaptive rule usable in the online HROM.





%===============================================================================
\section{Training manifold and local exactness systems}
%===============================================================================

Let the sampled latent manifold be
\[
        \Mtrain
        =
        \{\qMj\}_{j=1}^{\Nman},
        \qquad
        \qMj\in\R^{\Nlat}.
\]
In the current implementation,
\[
        \Nlat = 2,
        \qquad
        \qMj =
        \begin{bmatrix}
        q_{\mathrm{M},1}^{(j)} \\
        q_{\mathrm{M},2}^{(j)}
        \end{bmatrix}.
\]
The matrix used in the code is
\[
        \Qman =
        \begin{bmatrix}
        | & | &        & | \\
        \qM^{(1)} & \qM^{(2)} & \cdots & \qM^{(\Nman)} \\
        | & | &        & |
        \end{bmatrix}
        \in\R^{\Nlat\times\Nman}.
\]

For each sampled latent state $\qMj$, MAW--ECM constructs a local exactness
operator
\[
        \Uj \in \R^{\Ncand\times r_j},
\]
restricted to the current candidate support, and a right-hand side
\[
        \bj\in\R^{r_j}.
\]
The local adaptive weights $\wadapt^{(j)}\in\R^{\Ncand}$ must satisfy
\[
        \Uj^T\wadapt^{(j)}=\bj.
        \label{eq:local_exactness}
\]
When positivity is enforced, the additional condition is
\[
        \wadapt^{(j)}\geq 0.
        \label{eq:local_positivity}
\]

Equations~\eqref{eq:local_exactness}--\eqref{eq:local_positivity} define a
family of local cubature problems, one for each sampled manifold state. These
local problems are coupled only through the requirement that the support be
common to all states and, in the graph-regularized version, through the
smoothness penalty imposed on the weight fields.

%===============================================================================
\section{Fixed support and adaptive weights}
%===============================================================================

The defining structural choice of MAW--ECM is:
\[
        \text{fixed support, adaptive weights}.
\]
Thus,
\[
        \Zred
        \quad\text{is independent of }\qM,
\]
whereas
\[
        \wred = \wred(\qM).
\]
Equivalently, at the sampled manifold nodes,
\[
        \Wadapt_{\mathrm{red}}
        =
        \begin{bmatrix}
        | & | &        & | \\
        \wred(\qM^{(1)}) &
        \wred(\qM^{(2)}) &
        \cdots &
        \wred(\qM^{(\Nman)}) \\
        | & | &        & |
        \end{bmatrix}.
\]

This distinction is important. MAW--ECM is not an adaptive-support method. The
same integration points are evaluated online for all states. Therefore, the
finite element bookkeeping remains simple: the online code only needs the final
selected points and elements. What changes from state to state is the value of
the cubature weights.

A fully local cubature method would allow both
\[
        \Zred=\Zred(\qM),
        \qquad
        \wred=\wred(\qM).
\]
Such a strategy may produce even smaller local rules, but it would require
state-dependent point selection and would complicate online element
management. MAW--ECM deliberately avoids this route.

%===============================================================================
\section{Bivariate graph structure}
%===============================================================================

In the current implementation, the latent samples are assumed to lie on a
structured two-dimensional grid. The dimensions of the grid are denoted by
\[
        n_y,\; n_x,
        \qquad
        \Nman=n_y n_x .
\]
The two latent coordinate arrays are obtained by reshaping the rows of
$\Qman$:
\[
        Q_1 = \operatorname{reshape}
        \left(
        \{q_{\mathrm{M},1}^{(j)}\}_{j=1}^{\Nman},
        n_y,n_x
        \right),
\]
\[
        Q_2 = \operatorname{reshape}
        \left(
        \{q_{\mathrm{M},2}^{(j)}\}_{j=1}^{\Nman},
        n_y,n_x
        \right).
\]
A quadrilateral mesh is then built in the latent plane. The nodes of this mesh
are not physical finite element nodes; they are sampled reduced states. The
resulting graph is denoted by
\[
        \Ggraph=(\Mtrain,\Egraph),
\]
where $\Egraph$ contains the edges induced by the structured quadrilateral
connectivity.

The graph Laplacian, or graph stiffness matrix, is denoted by
\[
        \Kgraph\in\R^{\Nman\times\Nman}.
\]
For a scalar field $\varphi:\Mtrain\to\R$, collected in a vector
$\bm\varphi\in\R^{\Nman}$, the quantity
\[
        \bm\varphi^T\Kgraph\bm\varphi
\]
measures the roughness of the field over the sampled latent graph.

Since each retained integration point defines one scalar weight field over the
manifold, graph regularity of the adaptive weights can be measured by
\[
        \sum_{i=1}^{\Ncand}
        \bm w_i^T\Kgraph\bm w_i,
\]
where $\bm w_i\in\R^{\Nman}$ collects the values of the $i$-th integration-point
weight over all sampled states.

\paragraph{Specificity of the two-variable implementation.}
The use of a structured quadrilateral graph is specific to the bivariate
implementation. It relies on the possibility of reshaping the sampled manifold
states into a Cartesian-like array. For $\Nlat>2$, or for unstructured sampling
in any dimension, the same formulation remains valid if $\Kgraph$ is replaced
by a graph Laplacian constructed from a neighborhood relation in latent space,
for example using nearest neighbors, Delaunay connectivity, radial kernels, or
a manifold mesh.

%===============================================================================
\section{Global MAW--ECM algorithm}
%===============================================================================

The following pseudo-algorithm summarizes the logic implemented in
\texttt{MAW\_ECM\_genV1.m}. The individual routines invoked in the algorithm are
not expanded here; they are identified as separate sections to be developed
later.

%\input{alg_v1}
\begin{algorithm}[H]
\footnotesize
\caption{Offline construction of MAW--ECM with fixed support and adaptive weights}
\label{alg:maw_ecm_global}
\begin{algorithmic}[1]

\Require
Local integrand data $\{\Aint_j\}_{j=1}^{\Nman}$;
FE weights $\wFE$;
initial ECM support $\Zini$;
initial ECM weights $\wini$;
latent samples $\Qman$;
offline parameters.

\Ensure
Final support $\Zred$;
adaptive weight samples $\Wadapt_{\mathrm{red}}$;
regression model $\qM\mapsto\wred(\qM)$.

\vspace{0.2cm}

\State Build the constant scaffold
\[
        \Ufixed \gets \mathrm{orth}(\bone).
\]

\State Construct the latent graph and graph Laplacian
\[
        \Kgraph \gets
        \hyperref[sec:graph_laplacian]
        {\mathrm{GraphLaplacian}(\Qman)}.
\]

\State Build the weighted projector
\[
       \Qproj \gets \mathrm{WeightedOrthonormalBasis}(\Qall,\wFE).
\]

\State Construct the local exactness systems
\[
        \{\Uj,\bj\}_{j=1}^{\Nman}
        \gets
        \hyperref[sec:local_exactness_systems]{\mathrm{LocalBasisBuilder}
        (\Ufixed,\Zini,\{\Aint_j\},\wFE,\Qproj)}.
\]

\State Initialize the adaptive weights
\[
        \wadapt^{(j)} \gets \wini,
        \qquad
        j=1,\ldots,\Nman.
\]
Equivalently,
\[
        \Wadapt
        \gets
        [\wini,\ldots,\wini].
\]

\State Initialize active and eliminated sets
\[
        \Iact \gets \{1,\ldots,|\Zini|\},
        \qquad
        \Ielim \gets \emptyset.
\]

\While{stopping criterion is not satisfied}

    \State Compute aggregate candidate weights
    \[
        R_i = \sum_{j=1}^{\Nman} W_{ij}.
    \]

    \State Define positive candidate set
    \[
        \Ipos=\{i\in\Iact: R_i>0\}.
    \]

    \State Sort candidates in increasing order of $R_i$.

    \State Attempt a no-enforcement pruning step:
    \[
        (\Wadapt,\Iact,\Ielim,\mathrm{flag})
        \gets
        \hyperref[sec:no_enforcement_stage]{\mathrm{NoEnforcementPruning}
        (\{\Uj\},\{\bj\},\Wadapt,\Iact,\Ielim)}.
    \]

    \If{$\mathrm{flag}=\mathrm{success}$}

        \State Commit the accepted elimination.

    \Else

        \If{explicit positivity enforcement is enabled}

            \State Attempt graph-regularized active-set pruning:
            \[
                (\Wadapt,\Iact,\Ielim,\mathrm{flag})
                \gets
                \hyperref[sec:graph_regularized_stage]{\mathrm{GraphBasedPruning}
                (\{\Uj\},\{\bj\},\Wadapt,\Iact,\Ielim,\Kgraph)}.
            \]

        \Else

            \State Set $\mathrm{flag}=\mathrm{failure}$.

        \EndIf

        \If{$\mathrm{flag}=\mathrm{failure}$}
            \State Stop pruning.
        \EndIf

    \EndIf

\EndWhile

\State Extract the final support
\[
        \Zred \gets \Zini(\Iact).
\]

\State Extract the final adaptive weights
\[
        \Wadapt_{\mathrm{red}}
        \gets
        \Wadapt(\Iact,:).
\]

\State Regress the weight field
\[
        \qM \mapsto \wred(\qM)
\]
from the data
\[
        \left\{
        \qMj,\;
        \Wadapt_{\mathrm{red}}(:,j)
        \right\}_{j=1}^{\Nman}.
\]
\[
        \wred(\cdot)
        \gets
        \hyperref[sec:weight_regression]{\mathrm{WeightRegression}
        (\Qman,\Wadapt_{\mathrm{red}})}.
\]

\State Validate the adaptive cubature accuracy over the sampled manifold.

\end{algorithmic}
\end{algorithm}




Algorithm~\ref{alg:maw_ecm_global} summarizes the complete offline
construction of MAW--ECM. The construction of the local exactness systems is
detailed in Section~\ref{sec:local_exactness_systems}; the graph Laplacian
construction is described in Section~\ref{sec:graph_laplacian}; the three
pruning stages are presented in
Sections~\ref{sec:no_enforcement_stage},
\ref{sec:local_active_set_stage}, and
\ref{sec:graph_regularized_stage}; finally, the regression of the adaptive
weight fields is discussed in
Section~\ref{sec:weight_regression}.

%===============================================================================
\section{Interpretation of the algorithm}
%===============================================================================

The algorithm can be interpreted as a constrained compression of a family of
cubature rules. At the beginning, all sampled manifold states share the same
fixed ECM weights:
\[
        \wadapt^{(1)}
        =
        \wadapt^{(2)}
        =
        \cdots
        =
        \wadapt^{(\Nman)}
        =
        \wini.
\]
As pruning proceeds, the support is reduced, but the equality constraints are
maintained by allowing the columns of $\Wadapt$ to separate from one another.
Thus, the algorithm trades a reduction in the number of integration points for
an increase in the complexity of the weight map
\[
        \qM\mapsto\wred(\qM).
\]

This trade-off is the essential mechanism of MAW--ECM. A fixed ECM rule must
find one vector of weights that works globally over the whole manifold. MAW--ECM
instead constructs a family of related cubature rules:
\[
        \left\{
        \Zred,\wred(\qMj)
        \right\}_{j=1}^{\Nman},
\]
with common support but state-dependent weights.

The graph-regularized stage is introduced because the independent local
redistribution of weights may generate highly irregular weight fields. The graph
term couples neighboring latent states during the redistribution step and
discourages sharp jumps. However, it should not be interpreted as a guarantee of
global smoothness. The map $\qM\mapsto\wred(\qM)$ is the solution map of a
constrained optimization problem with active-set changes, and such maps are
generally only piecewise smooth.

%===============================================================================
\section{Specific and general components}
%===============================================================================

It is useful to separate the algorithm into components that are intrinsic to
MAW--ECM and components that are specific to the present bivariate
implementation.

\subsection{General MAW--ECM components}

The following components are independent of the latent dimension:
\begin{itemize}[leftmargin=1.2cm]
    \item the fixed-support/adaptive-weight formulation;
    \item the local exactness constraints $\Uj^T\wadapt^{(j)}=\bj$;
    \item the global pruning of integration points;
    \item the no-enforcement least-change redistribution;
    \item the positivity-constrained active-set redistribution;
    \item the final regression of $\wred(\qM)$;
    \item the validation of the reduced adaptive cubature rule.
\end{itemize}

\subsection{Components specific to the two-variable implementation}

The following components are specific to the present bivariate implementation:
\begin{itemize}[leftmargin=1.2cm]
    \item the assumption that $\Qman$ can be reshaped into an $n_y\times n_x$
    structured grid;
    \item the construction of a quadrilateral mesh in the latent plane;
    \item the finite-element-like graph Laplacian built on this quadrilateral
    mesh;
    \item the use of two-dimensional plotting and surface visualization of the
    adaptive weights;
    \item any hard-coded interpretation of the two latent variables as axial and
    shear-like macroscopic parameters.
\end{itemize}

\subsection{Extension to higher-dimensional latent manifolds}

For $\Nlat>2$, the algebraic structure of the method does not change. The local
constraints remain
\[
        \Uj^T\wadapt^{(j)}=\bj,
        \qquad
        j=1,\ldots,\Nman,
\]
and the adaptive weights are still stored as
\[
        \Wadapt\in\R^{\Ncand\times\Nman}.
\]
The only essential modification concerns the construction of $\Kgraph$.
Instead of a structured quadrilateral mesh in the $(q_{\mathrm{M},1},
q_{\mathrm{M},2})$ plane, one needs a graph on the sampled latent states in
$\R^{\Nlat}$. Once such a graph is available, the same graph energy
\[
        \sum_{i=1}^{\Ncand}\bm w_i^T\Kgraph\bm w_i
\]
can be used without changing the pruning formulation.

Thus, the two-dimensional implementation should be viewed as a convenient
realization of a more general graph-based MAW--ECM principle.

% %===============================================================================
% \section{Sections to be expanded next}
% %===============================================================================
%
% The following sections will be developed subsequently:
% \begin{enumerate}[leftmargin=1.2cm]
%     \item construction of the local basis matrices $\Uj$ and right-hand sides
%     $\bj$;
%     \item weighted projector and role of $\Qproj$;
%     \item construction of the latent graph and $\Kgraph$;
%     \item no-enforcement pruning;
%     \item local active-set positivity enforcement;
%     \item graph-regularized active-set pruning;
%     \item candidate ranking criteria;
%     \item RBF regression of the adaptive weights;
%     \item online evaluation and accuracy checks.
% \end{enumerate}


%===============================================================================
\section{Construction of the local exactness systems}
\label{sec:local_exactness_systems}
%===============================================================================

The first substantive operation in MAW--ECM is the construction of the local
linear systems that the adaptive weights must satisfy. These systems are the
manifold-dependent counterpart of the ECM exactness condition. In the classical
fixed-weight setting, a single cubature vector is sought so as to integrate a
global reduced basis of integrands. In MAW--ECM, by contrast, each sampled
manifold state carries its own local exactness system, while all systems share
the same candidate support.

Let the initial ECM rule be given by the support
\[
        \Zini = \{z_1,\ldots,z_{\Ncand}\}
        \subset \Zfull,
\]
and by the corresponding fixed weights
\[
        \wini \in \R^{\Ncand}.
\]
At this stage, $\Zini$ is not yet the final MAW--ECM support. It is the
candidate set from which the adaptive pruning algorithm will remove points.

For each sampled manifold state $\qMj$, let
\[
        \Aint_j \in \R^{\Nip\times m_j}
\]
denote the local integrand matrix associated with that state. The columns of
$\Aint_j$ contain the scalar fields over the full integration grid that one
wishes to integrate accurately at $\qMj$. Depending on the application, these
columns may represent stress-output integrands, residual integrands projected
onto a manifold tangent space, or a concatenation of several blocks.

The full finite element integral of these fields is
\[
        \Aint_j^T \wFE.
\]
However, the MAW--ECM pruning stage acts only on the initial candidate support
$\Zini$. Therefore, the local basis is restricted to the rows indexed by
$\Zini$:
\[
        \Uj
        =
        {\Ufull}_j(\Zini,:)
        \in \R^{\Ncand\times r_j},
\]
where ${\Ufull}_j$ denotes the orthonormalized local basis before restriction,
and $r_j$ is the number of local constraints retained at node $j$.

The local adaptive weights $\wadapt^{(j)}$ are required to satisfy
\[
        \Uj^T \wadapt^{(j)} = \bj,
        \qquad
        j=1,\ldots,\Nman.
        \label{eq:maw_local_system}
\]
The pair $(\Uj,\bj)$ is the fundamental local algebraic object used by all
subsequent pruning routines.

%-------------------------------------------------------------------------------
\subsection{Fixed scaffold}
%-------------------------------------------------------------------------------

A fixed scaffold is included in every local basis. Its role is to impose
constraints that must be preserved independently of the state-dependent
integrand content. In the present implementation, the most important scaffold
is the constant function,
\[
        \bone =
        (1,\ldots,1)^T\in\R^{\Nip}.
\]
After orthonormalization, this gives
\[
        \Ufixed = \operatorname{orth}(\bone).
\]
The corresponding constraint enforces preservation of the total measure carried
by the current cubature rule. In the candidate-support formulation, this means
that the adaptive weights should preserve the volume represented by the initial
ECM rule:
\[
        \bone_{\Zini}^T \wadapt^{(j)}
        =
        \bone_{\Zini}^T \wini,
        \qquad
        j=1,\ldots,\Nman.
\]
Thus, even if the state-dependent integrand block becomes small or degenerate,
the adaptive rule is prevented from collapsing to the trivial zero solution.

More general scaffolds are possible. For instance, one may include additional
low-order functions, elastic modes, or previously identified integration
constraints. In that case,
\[
        \Ufixed
        =
        \operatorname{orth}
        \left(
        [\Uel\;\;\bone]
        \right),
\]
where $\Uel$ denotes a user-defined fixed block. The present bivariate
implementation uses only the constant scaffold.

%-------------------------------------------------------------------------------
\subsection{Weighted projector}
%-------------------------------------------------------------------------------

The routine also constructs a weighted projector from a precomputed basis
$\Qall$. This object is used to project or filter the local integrand matrices
before building the local exactness bases.

Let
\[
        \Qall \in \R^{\Nip\times n_Q}
\]
denote the input basis. The code forms
\[
        \Qproj
        =
        \operatorname{orth}
        \left(
        \diag(\wFE)^{-1/2}\Qall
        \right).
\]
Equivalently, if
\[
        \bm S_{\mathrm{FE}}
        =
        \diag(\sqrt{w^{\mathrm{FE}}_1},\ldots,\sqrt{w^{\mathrm{FE}}_{\Nip}}),
\]
then
\[
        \Qproj
        =
        \operatorname{orth}
        \left(
        \bm S_{\mathrm{FE}}^{-1}\Qall
        \right).
\]
This transformation is intended to work with a basis compatible with the
finite-element weighted inner product. In implementation terms, this is the
operation
\[
        \Qproj
        \leftarrow
        \operatorname{SVDT}
        \left(
        \Qall ./ \sqrt{\wFE}
        \right).
\]

When the projection option is active, the local integrand matrix is replaced by
\[
        \Aint_j^{\mathrm{proj}}
        =
        \Qproj(\Qproj^T\Aint_j).
\]
This projected content is then used to construct the local basis. The purpose
is not to change the cubature philosophy, but to avoid imposing local
constraints associated with high-frequency or poorly represented components of
the integrand snapshots.

%-------------------------------------------------------------------------------
\subsection{Local basis at a sampled manifold state}
%-------------------------------------------------------------------------------

For each sampled state $j$, one starts from the local integrand block
$\Aint_j$. If the projection option is active, the projected matrix
$\Aint_j^{\mathrm{proj}}$ is used instead. For compactness, denote the matrix
actually used by
\[
        \widehat{\Aint}_j.
\]
Thus,
\[
        \widehat{\Aint}_j =
        \begin{cases}
        \Qproj(\Qproj^T\Aint_j),
        & \text{if projection is active},\\[2mm]
        \Aint_j,
        & \text{otherwise}.
        \end{cases}
\]

The local full-grid basis is then obtained by orthonormalizing the
concatenation
\[
        {\Ufull}_j
        =
        \operatorname{orth}
        \left(
        [\Ufixed\;\;\widehat{\Aint}_j]
        \right).
\]
In the implementation this orthonormalization is performed by \texttt{SVDT}.
Only the rows associated with the initial ECM support are retained:
\[
        \Uj = {\Ufull}_j(\Zini,:).
\]

The local right-hand side depends on whether projection is active. Without
projection, the natural target is the full FE integral:
\[
        \bj
        =
        {\Ufull}_j^T\wFE.
        \label{eq:b_full_fe}
\]
With projection, the implementation instead uses the integral induced by the
current fixed ECM rule:
\[
        \bj
        =
        \Uj^T\wini.
        \label{eq:b_projected_ecm}
\]
This distinction is important. In the projected case, the local system is made
consistent with the current candidate rule. The pruning problem then asks
whether the same projected constraints can still be satisfied after removing
additional points from $\Zini$.

%-------------------------------------------------------------------------------
\subsection{Initialization of the adaptive weights}
%-------------------------------------------------------------------------------

Once the local systems have been constructed, the adaptive weight matrix is
initialized by copying the fixed ECM weights at every sampled manifold state:
\[
        \wadapt^{(j)} = \wini,
        \qquad
        j=1,\ldots,\Nman.
\]
Equivalently,
\[
        \Wadapt
        =
        [\wini,\wini,\ldots,\wini]
        \in\R^{\Ncand\times\Nman}.
\]
Thus, the initial MAW--ECM rule is exactly the fixed ECM rule replicated over
the sampled manifold. Adaptivity is introduced only through the subsequent
pruning and redistribution steps.

This initialization has two useful consequences. First, feasibility is inherited
from the initial rule, up to the consistency of the constructed local systems.
Second, the pruning algorithm starts from a state in which the weight field is
constant over the manifold:
\[
        \wadapt(\qM^{(1)})
        =
        \cdots
        =
        \wadapt(\qM^{(\Nman)}).
\]
All spatial variability of the adaptive weights is therefore generated by the
need to preserve the local exactness constraints after points are removed.

%-------------------------------------------------------------------------------
\subsection{Function-level interpretation}
%-------------------------------------------------------------------------------

The construction described in this section is implemented in
\[
        \texttt{MAWecmBasisMATRIX\_LOCAL\_noverlap.m}.
\]
Its role is purely preparatory: it does not prune integration points and does
not regress the adaptive weights. It builds the local systems
\[
        \{(\Uj,\bj)\}_{j=1}^{\Nman}
\]
and initializes $\Wadapt$.

In pseudo-code, the function performs the following operations.

\begin{algorithm}[H]
\caption{Construction of local exactness systems}
\label{alg:local_exactness_systems}
\begin{algorithmic}[1]

\Require
Fixed scaffold $\Ufixed$;
initial support $\Zini$;
local integrand matrices $\{\Aint_j\}_{j=1}^{\Nman}$;
initial weights $\wini$;
FE weights $\wFE$;
optional projector $\Qproj$.

\Ensure
Local operators $\{\Uj\}$;
right-hand sides $\{\bj\}$;
initial adaptive weights $\Wadapt$.

\For{$j=1,\ldots,\Nman$}

    \If{projection through $\Qproj$ is active}
        \State Compute
        \[
            \widehat{\Aint}_j
            =
            \Qproj(\Qproj^T\Aint_j).
        \]
    \Else
        \State Set
        \[
            \widehat{\Aint}_j=\Aint_j.
        \]
    \EndIf

    \State Build the full local basis
    \[
        {\Ufull}_j
        =
        \operatorname{orth}
        \left(
        [\Ufixed\;\;\widehat{\Aint}_j]
        \right).
    \]

    \State Restrict the basis to the current candidate support
    \[
        \Uj={\Ufull}_j(\Zini,:).
    \]

    \If{projection through $\Qproj$ is active}
        \State Define the right-hand side by the current ECM rule
        \[
            \bj=\Uj^T\wini.
        \]
    \Else
        \State Define the right-hand side by the full FE quadrature
        \[
            \bj={\Ufull}_j^T\wFE.
        \]
    \EndIf

    \State Initialize
    \[
        \wadapt^{(j)}=\wini.
    \]

\EndFor

\State Assemble
\[
        \Wadapt =
        [\wadapt^{(1)},\ldots,\wadapt^{(\Nman)}].
\]

\end{algorithmic}
\end{algorithm}

%-------------------------------------------------------------------------------
\subsection{Specificity and possible generalizations}
%-------------------------------------------------------------------------------

The construction of the local exactness systems is not specific to a
two-dimensional latent manifold. The only required input is a collection of
sampled states indexed by $j=1,\ldots,\Nman$. Therefore, the equations
\[
        \Uj^T\wadapt^{(j)}=\bj
\]
remain unchanged for any latent dimension $\Nlat$.

What is specific to the present implementation is the organization of the
samples and the later use of a structured bivariate graph. The local systems
themselves do not know whether the manifold is one-dimensional, two-dimensional,
or higher-dimensional. They only require a consistent ordering between:
\[
        \qMj,\quad
        \Aint_j,\quad
        \Uj,\quad
        \bj,\quad
        \wadapt^{(j)}.
\]

This observation is important for future extensions. In higher-dimensional or
unstructured latent spaces, the present section would remain essentially
unchanged. Only the graph construction and visualization tools would need to be
modified.


%===============================================================================
\section{Graph Laplacian on the sampled latent manifold}
\label{sec:graph_laplacian}
%===============================================================================

The graph-based version of MAW--ECM requires a discrete notion of smoothness for
scalar fields defined over the sampled latent manifold. In the present setting,
these scalar fields are the adaptive cubature weights associated with the
retained integration points. For a fixed integration point $i$, the values
\[
        W_{i1},W_{i2},\ldots,W_{i\Nman}
\]
define a scalar field over the sampled latent states
\[
        \{\qMj\}_{j=1}^{\Nman}.
\]
The purpose of the graph Laplacian is to measure and penalize nonsmooth
variations of such fields across neighboring latent states.

The graph construction is independent of the physical finite element mesh. It
is built in the latent space, not in the spatial domain. The physical mesh
defines the integration points; the latent graph defines neighborhood relations
between sampled reduced states.

%-------------------------------------------------------------------------------
\subsection{Structured bivariate graph}
%-------------------------------------------------------------------------------

In the current implementation, the latent dimension is two:
\[
        \Nlat = 2,
        \qquad
        \qM =
        \begin{bmatrix}
        q_{\mathrm{M},1}\\
        q_{\mathrm{M},2}
        \end{bmatrix}.
\]
The sampled states are assumed to be ordered as a structured grid. Let
\[
        n_y n_x = \Nman
\]
be the number of samples in the two grid directions. The two rows of $\Qman$
are reshaped into two arrays
\[
        \Qone =
        \operatorname{reshape}
        \left(
        q_{\mathrm{M},1}^{(1)},\ldots,q_{\mathrm{M},1}^{(\Nman)};
        n_y,n_x
        \right),
\]
and
\[
        \Qtwo =
        \operatorname{reshape}
        \left(
        q_{\mathrm{M},2}^{(1)},\ldots,q_{\mathrm{M},2}^{(\Nman)};
        n_y,n_x
        \right).
\]
Each entry $(a,b)$ of these arrays defines one node of the latent graph:
\[
        \bm x^{\mathrm{G}}_{ab}
        =
        \begin{bmatrix}
        \Qone(a,b)\\
        \Qtwo(a,b)
        \end{bmatrix}.
\]
A structured quadrilateral mesh is then built over these nodes. Its
connectivity is denoted by
\[
        \mathcal{T}_{\mathrm{G}}.
\]
Again, this is not a finite element mesh in the physical domain. It is an
auxiliary mesh over the sampled latent states.

%-------------------------------------------------------------------------------
\subsection{Finite-element-like graph stiffness}
%-------------------------------------------------------------------------------

The graph Laplacian $\Kgraph$ is assembled as a scalar diffusion stiffness
matrix on the auxiliary latent mesh. Formally, if $\varphi$ is a scalar field
defined over the latent manifold, one may view its nodal values as
\[
        \bm\varphi =
        \begin{bmatrix}
        \varphi(\qM^{(1)}) &
        \cdots &
        \varphi(\qM^{(\Nman)})
        \end{bmatrix}^T .
\]
The discrete graph energy is then
\[
        \mathcal{E}_{\mathrm{G}}(\bm\varphi)
        =
        \bm\varphi^T\Kgraph\bm\varphi.
        \label{eq:graph_energy_scalar}
\]
This quantity is the discrete counterpart of a Dirichlet energy in latent
space,
\[
        \int_{\mathcal{M}}
        \nabla_{\qM}\varphi\cdot\nabla_{\qM}\varphi
        \,d\qM.
\]
Thus, $\mathcal{E}_{\mathrm{G}}$ is small for slowly varying fields and large
for fields with strong jumps between neighboring latent states.

In MAW--ECM, the scalar fields of interest are the rows of the adaptive weight
matrix. Let
\[
        \bm w_i =
        \begin{bmatrix}
        W_{i1} & W_{i2} & \cdots & W_{i\Nman}
        \end{bmatrix}^T
        \in\R^{\Nman}
\]
be the weight field associated with integration point $i$. The total graph
roughness of the adaptive rule is measured by
\[
        \mathcal{S}_{\mathrm{G}}(\Wadapt)
        =
        \sum_{i=1}^{\Ncand}
        \bm w_i^T\Kgraph\bm w_i.
        \label{eq:graph_energy_weights}
\]
Equivalently, since the rows of $\Wadapt$ are the weight fields,
\[
        \mathcal{S}_{\mathrm{G}}(\Wadapt)
        =
        \trace\left(
        \Wadapt\Kgraph\Wadapt^T
        \right).
        \label{eq:graph_energy_trace}
\]

%-------------------------------------------------------------------------------
\subsection{Incremental smoothing}
%-------------------------------------------------------------------------------

In the pruning step, it is usually preferable to regularize the change in the
weights rather than the absolute value of the weights. Let
\[
        \DeltaW = \Wnew-\Wold
\]
be the redistribution induced by a tentative point removal. The graph
regularization term used in the update can then be written as
\[
        \mathcal{S}_{\mathrm{G}}(\DeltaW)
        =
        \trace\left(
        \DeltaW\Kgraph\DeltaW^T
        \right).
        \label{eq:incremental_graph_energy}
\]
This choice has a clear interpretation. The previous feasible weights
$\Wold$ are taken as a reference state. The update is allowed to change them,
but abrupt changes over neighboring latent states are penalized.

A representative graph-regularized least-change objective is therefore
\[
        \frac{1}{2}\|\Wnew-\Wold\|_F^2
        +
        \frac{\alphaG}{2}
        \trace\left[
        (\Wnew-\Wold)\Kgraph(\Wnew-\Wold)^T
        \right],
        \label{eq:incremental_graph_objective}
\]
where $\alphaG\geq 0$ controls the strength of the graph regularization.
For $\alphaG=0$, one recovers a purely local least-change redistribution. As
$\alphaG$ increases, rough variations of the redistribution over the latent
graph become increasingly penalized.

%-------------------------------------------------------------------------------
\subsection{Stacked-vector form}
%-------------------------------------------------------------------------------

For the graph-regularized active-set update, it is useful to write
\eqref{eq:incremental_graph_objective} in vector form. Let
\[
        \zvec = \vecop(\Wnew),
        \qquad
        \zold = \vecop(\Wold),
\]
where the columns of the weight matrix are stacked. Then
\[
        \|\Wnew-\Wold\|_F^2
        =
        \|\zvec-\zold\|_2^2.
\]
Moreover,
\[
        \trace\left[
        (\Wnew-\Wold)\Kgraph(\Wnew-\Wold)^T
        \right]
        =
        (\zvec-\zold)^T
        (\Kgraph\otimes\bm I_{\Ncand})
        (\zvec-\zold).
\]
Consequently, the incremental graph-regularized objective becomes
\[
        \frac{1}{2}
        (\zvec-\zold)^T
        \left[
        \bm I + \alphaG(\Kgraph\otimes\bm I_{\Ncand})
        \right]
        (\zvec-\zold).
        \label{eq:stacked_graph_objective}
\]
Define
\[
        \Hgraph
        :=
        \bm I + \alphaG(\Kgraph\otimes\bm I_{\Ncand}).
        \label{eq:Hgraph_definition}
\]
Then the objective is simply
\[
        \frac{1}{2}
        (\zvec-\zold)^T\Hgraph(\zvec-\zold).
\]

The structure of the Kronecker product is important. The Laplacian
$\Kgraph$ acts across manifold samples, whereas $\bm I_{\Ncand}$ acts on the
integration-point index. Therefore, the graph regularization does not mix
different integration points. Each weight field is smoothed independently over
the manifold graph.

%-------------------------------------------------------------------------------
\subsection{Boundary-weighted graph regularization}
%-------------------------------------------------------------------------------

The implementation also allows the graph Laplacian to be modified near the
boundary of the sampled latent domain. This option increases the regularization
strength close to the boundary, where interpolation and extrapolation effects
are often more delicate.

Let $d_j$ denote the grid-index distance from node $j$ to the closest boundary
of the structured latent grid. A nodal amplification factor is defined as
\[
        \rho_j
        =
        1+\lambda_{\mathrm{B}}
        \exp\left(-\frac{d_j}{\ell_{\mathrm{B}}}\right),
        \qquad
        \rho_j\geq 1.
\]
Here $\lambda_{\mathrm{B}}$ controls the strength of the boundary
amplification, and $\ell_{\mathrm{B}}$ controls its decay length.

For an edge connecting nodes $a$ and $b$, the edge amplification is taken as
\[
        \rho_{ab}
        =
        \frac{1}{2}(\rho_a+\rho_b).
\]
The original graph edge weight is then multiplied by $\rho_{ab}$, and the
Laplacian is rebuilt from the modified edge weights. The resulting matrix is a
boundary-weighted graph Laplacian. Its purpose is not to change the exactness
constraints, but only to alter the smoothness metric used during graph-based
redistribution.

%-------------------------------------------------------------------------------
\subsection{Specificity and extension to higher dimensions}
%-------------------------------------------------------------------------------

The construction based on $\Qone$, $\Qtwo$ and quadrilateral connectivity is
specific to the bivariate implementation. It requires that the latent samples
can be reshaped into a structured two-dimensional grid.

The graph-regularized MAW--ECM formulation itself does not require this. For a
general latent dimension $\Nlat$, one only needs a graph
\[
        \Ggraph=(\Mtrain,\Egraph)
\]
whose nodes are the sampled latent states. The graph Laplacian can then be
constructed from any suitable neighborhood relation. Possible choices include:
\[
        k\text{-nearest-neighbor graphs},\qquad
        \varepsilon\text{-graphs},\qquad
        \text{radial-kernel graphs},\qquad
        \text{simplicial or manifold meshes}.
\]
Once $\Kgraph$ is available, the same graph energy
\[
        \trace(\Wadapt\Kgraph\Wadapt^T)
\]
and the same stacked operator
\[
        \bm I+\alphaG(\Kgraph\otimes\bm I)
\]
apply without modification.

Thus, the bivariate graph used in the present implementation should be viewed
as one convenient realization of a more general graph-based smoothing principle
for manifold-adaptive cubature weights.

%-------------------------------------------------------------------------------
\subsection{Function-level interpretation}
%-------------------------------------------------------------------------------

The construction of the bivariate graph is split into small auxiliary routines.
The function
\[
        \texttt{quadMeshFromXY.m}
\]
constructs the nodal coordinate array and quadrilateral connectivity from the
reshaped latent-coordinate arrays $\Qone$ and $\Qtwo$. The function
\[
        \texttt{ComputeK\_scl.m}
\]
assembles the scalar finite-element-like graph stiffness matrix $\Kgraph$ on
this auxiliary latent mesh. If boundary amplification is activated,
\[
        \texttt{BoundaryWeightedGraphLaplacian.m}
\]
modifies the edge weights and rebuilds the Laplacian.

The resulting matrix $\Kgraph$ is then passed to the graph-based pruning
routine, where it enters the regularized redistribution of the adaptive
weights.


%===============================================================================
\section{Pruning strategy and candidate elimination}
\label{sec:pruning_strategy}
%===============================================================================

Once the local exactness systems and the graph operator have been constructed,
MAW--ECM enters the pruning stage. The objective is to reduce the initial
candidate support $\Zini$ while preserving, at every sampled manifold state,
the local cubature constraints
\[
        \Uj^T\wadapt^{(j)}=\bj,
        \qquad
        j=1,\ldots,\Nman.
\]
The reduction is performed by eliminating integration points one at a time.
Each elimination removes one row of the adaptive weight matrix
\[
        \Wadapt\in\R^{\Ncand\times\Nman}.
\]
Thus, pruning is global with respect to the support: if a point is eliminated,
it is eliminated for all manifold states.

The adaptive character of the method enters through the redistribution of the
remaining weights. After a point is removed, each column of $\Wadapt$ is allowed
to change so that the corresponding local exactness system remains satisfied.

%-------------------------------------------------------------------------------
\subsection{Active and eliminated sets}
%-------------------------------------------------------------------------------

Let
\[
        \Iact \subset \{1,\ldots,\Ncand\}
\]
denote the set of currently active candidate indices, and let
\[
        \Ielim
        =
        \{1,\ldots,\Ncand\}\setminus\Iact
\]
denote the eliminated indices. At the beginning of the pruning stage,
\[
        \Iact = \{1,\ldots,\Ncand\},
        \qquad
        \Ielim=\emptyset.
\]

If $p\in\Iact$ is selected for elimination, the updated active set is
\[
        \Iact^{+}
        =
        \Iact\setminus\{p\},
\]
and the eliminated set becomes
\[
        \Ielim^{+}
        =
        \Ielim\cup\{p\}.
\]
In terms of adaptive weights, the elimination means
\[
        W_{pj}^{+}=0,
        \qquad
        j=1,\ldots,\Nman.
\]
The remaining rows of $\Wadapt$ must then be recomputed.

%-------------------------------------------------------------------------------
\subsection{Candidate ranking by aggregate weight}
%-------------------------------------------------------------------------------

The first candidate ranking used in the algorithm is based on the aggregate
contribution of each integration point over the sampled manifold. For a current
weight matrix $\Wadapt$, define
\[
        R_i
        =
        \sum_{j=1}^{\Nman} W_{ij},
        \qquad
        i=1,\ldots,\Ncand.
        \label{eq:aggregate_weight}
\]
Only candidates with positive aggregate contribution are considered:
\[
        \Ipos
        =
        \{i\in\Iact : R_i>0\}.
\]
The candidates in $\Ipos$ are sorted in increasing order of $R_i$. Hence, the
algorithm first attempts to remove the integration points that carry the
smallest total weight over the manifold.

This criterion is intentionally simple. It is the natural extension of the
greedy intuition used in fixed-weight cubature: points with small cumulative
weight are expected to be less important. In MAW--ECM, however, a small value of
$R_i$ does not guarantee that the point can be safely removed. A point may have
small total weight and still be essential for satisfying one or several local
exactness systems.

%-------------------------------------------------------------------------------
\subsection{Single-point elimination trial}
%-------------------------------------------------------------------------------

Consider a candidate point $p\in\Ipos$. A pruning trial consists of imposing
\[
        W_{pj}^{+}=0,
        \qquad
        j=1,\ldots,\Nman,
        \label{eq:global_prune_constraint}
\]
and seeking updated weights $\Wnew$ such that
\[
        \Uj^T{\wnew}^{(j)}=\bj,
        \qquad
        j=1,\ldots,\Nman.
        \label{eq:post_prune_exactness}
\]
If positivity is enforced, one also requires
\[
        {\wnew}^{(j)}\geq 0,
        \qquad
        j=1,\ldots,\Nman.
        \label{eq:post_prune_positivity}
\]

A candidate is accepted only if the redistribution problem is feasible according
to the active stage of the algorithm. Otherwise, the point is rejected and the
next candidate in the ordered list is tested.

The important point is that at most one integration point is committed per
successful pruning pass. This conservative strategy makes the evolution of the
support easier to monitor and reduces the risk of destroying feasibility by an
overly aggressive simultaneous elimination.

%-------------------------------------------------------------------------------
\subsection{Stopping criterion}
%-------------------------------------------------------------------------------

The pruning loop terminates when either the target number of points has been
reached or no further feasible elimination can be found. A natural lower bound
is dictated by the number of local constraints. If
\[
        r_j = \dim(\bj),
\]
then a necessary condition for local feasibility is that the number of active
points should not fall below the number of independent constraints at the most
demanding node. Hence one defines
\[
        r_{\max}
        =
        \max_{1\leq j\leq\Nman} r_j.
\]
In practice, an additional user-prescribed maximum or target number of ECM
points may also be imposed. Denoting this value by $n_{\mathrm{target}}$, the
effective stopping threshold can be written as
\[
        n_{\mathrm{stop}}
        =
        \max(r_{\max},n_{\mathrm{target}}).
\]
The pruning loop is not allowed to reduce the support below this threshold.

Even above this threshold, pruning may stop earlier if no candidate produces a
feasible redistribution. This is a stronger and more meaningful stopping
criterion than the mere count of active points.

%-------------------------------------------------------------------------------
\subsection{Two-stage pruning logic}
%-------------------------------------------------------------------------------

The implementation distinguishes two main redistribution stages.

The first stage attempts a least-change update without explicit positivity
enforcement. This stage is computationally inexpensive and usually produces
smooth changes in the adaptive weights. Symbolically, for a candidate point
$p$, it resembles
\[
        \min_{\Wnew}
        \frac{1}{2}\|\Wnew-\Wold\|_F^2,
\]
subject to
\[
        \Uj^T{\wnew}^{(j)}=\bj,
        \qquad
        W_{pj}^{+}=0,
        \qquad
        j=1,\ldots,\Nman.
\]
No inequality constraint is imposed at this stage.

If this no-enforcement stage can no longer eliminate points, the algorithm may
switch to an explicit positivity stage. In that case, the redistribution also
enforces
\[
        {\wnew}^{(j)}\geq 0.
\]
This positivity-enforcing stage can be performed either independently at each
manifold node, or through the graph-regularized coupled formulation described
later.

The transition between these stages is important. The no-enforcement stage
usually preserves a relatively regular weight field. The explicit positivity
stage is more robust from the cubature viewpoint, but it introduces active-set
changes, and these active-set changes are a primary source of nonsmoothness in
the mapping
\[
        \qM\mapsto\wred(\qM).
\]

%-------------------------------------------------------------------------------
\subsection{Pseudo-algorithm for the pruning loop}
%-------------------------------------------------------------------------------

The pruning loop can be summarized as follows.

\begin{algorithm}[H]
\caption{MAW--ECM pruning loop}
\label{alg:maw_pruning_loop}
\begin{algorithmic}[1]

\Require
Local systems $\{(\Uj,\bj)\}_{j=1}^{\Nman}$;
initial adaptive weights $\Wadapt$;
initial support $\Zini$;
graph Laplacian $\Kgraph$;
stopping threshold $n_{\mathrm{stop}}$.

\Ensure
Reduced support $\Zred$;
adaptive weights on the reduced support.

\State Initialize
\[
        \Iact\gets\{1,\ldots,|\Zini|\},
        \qquad
        \Ielim\gets\emptyset.
\]

\While{$|\Iact|>n_{\mathrm{stop}}$}

    \State Compute aggregate weights
    \[
        R_i=\sum_{j=1}^{\Nman}W_{ij}.
    \]

    \State Define the candidate list
    \[
        \Ipos=\{i\in\Iact:R_i>0\}.
    \]

    \State Sort $\Ipos$ in increasing order of $R_i$.

    \State Attempt no-enforcement pruning over the ordered candidates.

    \If{a feasible candidate is found}
        \State Commit the elimination:
        \[
            \Iact\gets\Iact\setminus\{p\},
            \qquad
            \Ielim\gets\Ielim\cup\{p\}.
        \]
        \State Update $\Wadapt$.
    \Else
        \If{positivity enforcement is disabled}
            \State Stop.
        \Else
            \If{graph-based redistribution is enabled}
                \State Attempt graph-regularized active-set pruning.
            \Else
                \State Attempt local active-set pruning independently at each node.
            \EndIf

            \If{a feasible candidate is found}
                \State Commit the elimination and update $\Wadapt$.
            \Else
                \State Stop.
            \EndIf
        \EndIf
    \EndIf

\EndWhile

\State Set
\[
        \Zred=\Zini(\Iact).
\]

\State Keep
\[
        \Wadapt_{\mathrm{red}}=\Wadapt(\Iact,:).
\]

\end{algorithmic}
\end{algorithm}

%-------------------------------------------------------------------------------
\subsection{Interpretation}
%-------------------------------------------------------------------------------

The pruning process should not be viewed as simply deleting points with small
weights. Each deletion changes the feasible set of all local cubature systems.
A point can be removed only if the remaining support can still reproduce the
local constraints at every sampled manifold state.

Thus, MAW--ECM performs a constrained sparsification of a family of cubature
rules. The support is sparse and common, but the weights are allowed to move
over the manifold. This additional freedom is precisely what allows the support
to be smaller than in a fixed-weight ECM rule.

At the same time, this freedom introduces a new difficulty. The adaptive weight
field is the result of repeated constrained redistributions. When positivity is
enforced, changes of active set may produce kinks, jumps, or localized spikes.
The graph-based stage is introduced to moderate these effects, but the
underlying constrained nature of the problem remains.

%-------------------------------------------------------------------------------
\subsection{Function-level interpretation}
%-------------------------------------------------------------------------------

The global pruning loop is orchestrated by
\[
        \texttt{MAW\_ECM\_genV1.m}.
\]
The no-enforcement stage is implemented in
\[
        \texttt{Loop\_MAWecmNOENF\_cl.m}.
\]
The explicit positivity-enforcing local stage is implemented in
\[
        \texttt{Loop\_MAWecmNOENFe.m}.
\]
The graph-based positivity-enforcing stage is implemented in
\[
        \texttt{GraphBased\_MAW\_ECM\_2v.m},
\]
which in turn calls
\[
        \texttt{prune\_step\_optionB.m}.
\]

The detailed mathematical formulation of these redistribution stages is not
included in the present section. It will be developed in the following sections,
starting from the no-enforcement least-change update and then moving to the
active-set and graph-regularized formulations.


%===============================================================================
\section{Stage 1: no-enforcement least-change redistribution}
\label{sec:no_enforcement_stage}
%===============================================================================

The first redistribution stage of MAW--ECM attempts to remove integration
points without explicitly enforcing positivity of the adaptive weights. This
stage is implemented in
\[
        \texttt{Loop\_MAWecmNOENF\_cl.m}.
\]
Its role is to perform inexpensive pruning steps while preserving the local
exactness constraints with minimal perturbation of the current adaptive
weights.

This stage is important for two reasons. First, it is computationally cheaper
than the active-set positivity-enforcing stages. Second, and more importantly,
it tends to preserve smoothness of the adaptive weight field. Since no weights
are explicitly clamped to zero, the redistribution remains governed only by the
equality constraints, and the resulting changes in the weights are usually
relatively gentle over the latent manifold.

The no-enforcement stage should therefore be interpreted as the preferred
redistribution regime. The explicit positivity stages are activated only when
the unconstrained redistribution is no longer able to reduce the support.

%-------------------------------------------------------------------------------
\subsection{Local least-change redistribution}
%-------------------------------------------------------------------------------

Suppose that the current candidate support contains $\Ncand$ active points and
that a candidate point $p$ is selected for elimination. The pruning condition
is
\[
        W_{pj}^{+}=0,
        \qquad
        j=1,\ldots,\Nman.
\]
The remaining adaptive weights must then be redistributed so as to preserve the
local exactness systems
\[
        {\Uj}^T{\wnew}^{(j)}=\bj,
        \qquad
        j=1,\ldots,\Nman.
\]
At the no-enforcement stage, the redistribution is performed independently at
each manifold sample.

Let
\[
        {\wold}^{(j)}
\]
denote the current local adaptive weights at node $j$. The updated local
weights are obtained by solving a least-change problem of the form
\[
        \min_{{\wnew}^{(j)}}
        \frac{1}{2}
        \|
        {\wnew}^{(j)}-{\wold}^{(j)}
        \|_2^2,
        \label{eq:no_enforcement_objective}
\]
subject to
\[
        \Uj^T{\wnew}^{(j)}=\bj,
        \label{eq:no_enforcement_exactness}
\]
and
\[
        w_{p,\mathrm{new}}^{(j)}=0.
        \label{eq:no_enforcement_elimination}
\]
No positivity condition is imposed:
\[
        {\wnew}^{(j)}
        \not\geq 0
        \quad\text{necessarily}.
\]

Thus, the update seeks the feasible redistribution closest, in Euclidean norm,
to the previous local weights.

%-------------------------------------------------------------------------------
\subsection{Reduced local system after elimination}
%-------------------------------------------------------------------------------

Let
\[
        \Irem
        =
        \Iact\setminus\{p\}
\]
denote the remaining active rows after tentative elimination of point $p$.
Define the reduced local basis
\[
        \widetilde{\Uj}
        =
        \Uj(\Irem,:)
        \in\R^{(\Ncand-1)\times r_j}.
\]
Similarly, define the reduced local unknown
\[
        \widetilde{\wnew}^{(j)}
        =
        {\wnew}^{(j)}(\Irem).
\]
Since the eliminated weight is fixed to zero, the local exactness condition
reduces to
\[
        \widetilde{\Uj}^T
        \widetilde{\wnew}^{(j)}
        =
        \bj.
        \label{eq:reduced_exactness_system}
\]
The least-change problem therefore becomes
\[
        \min_{\widetilde{\wnew}^{(j)}}
        \frac{1}{2}
        \|
        \widetilde{\wnew}^{(j)}
        -
        \widetilde{\wold}^{(j)}
        \|_2^2,
\]
subject to
\[
        \widetilde{\Uj}^T
        \widetilde{\wnew}^{(j)}
        =
        \bj.
        \label{eq:reduced_least_change_problem}
\]

%-------------------------------------------------------------------------------
\subsection{Closed-form least-change update}
%-------------------------------------------------------------------------------

The problem
\eqref{eq:reduced_least_change_problem} is a standard equality-constrained
quadratic minimization problem. Introducing Lagrange multipliers
\[
        \bm\lambda_j\in\R^{r_j},
\]
the local Lagrangian reads
\[
        \mathcal{L}
        =
        \frac{1}{2}
        \|
        \widetilde{\wnew}^{(j)}
        -
        \widetilde{\wold}^{(j)}
        \|_2^2
        +
        \bm\lambda_j^T
        \left(
        \widetilde{\Uj}^T
        \widetilde{\wnew}^{(j)}
        -
        \bj
        \right).
\]
Stationarity with respect to
$\widetilde{\wnew}^{(j)}$ gives
\[
        \widetilde{\wnew}^{(j)}
        -
        \widetilde{\wold}^{(j)}
        +
        \widetilde{\Uj}\bm\lambda_j
        =
        0.
\]
Hence,
\[
        \widetilde{\wnew}^{(j)}
        =
        \widetilde{\wold}^{(j)}
        -
        \widetilde{\Uj}\bm\lambda_j.
        \label{eq:no_enforcement_update}
\]
Substituting into the constraint yields
\[
        \widetilde{\Uj}^T
        \widetilde{\Uj}
        \bm\lambda_j
        =
        \widetilde{\Uj}^T
        \widetilde{\wold}^{(j)}
        -
        \bj.
        \label{eq:no_enforcement_multiplier}
\]
If the Gram matrix
\[
        \widetilde{\Uj}^T\widetilde{\Uj}
\]
is invertible, one obtains
\[
        \bm\lambda_j
        =
        \left(
        \widetilde{\Uj}^T\widetilde{\Uj}
        \right)^{-1}
        \left(
        \widetilde{\Uj}^T
        \widetilde{\wold}^{(j)}
        -
        \bj
        \right).
\]
The least-change correction therefore becomes
\[
        \widetilde{\wnew}^{(j)}
        =
        \widetilde{\wold}^{(j)}
        -
        \widetilde{\Uj}
        \left(
        \widetilde{\Uj}^T\widetilde{\Uj}
        \right)^{-1}
        \left(
        \widetilde{\Uj}^T
        \widetilde{\wold}^{(j)}
        -
        \bj
        \right).
        \label{eq:orthogonal_projection_update}
\]
This expression shows that the update is an orthogonal projection of the
previous weights onto the affine subspace defined by the reduced exactness
constraints.

%-------------------------------------------------------------------------------
\subsection{Interpretation as a minimum redistribution}
%-------------------------------------------------------------------------------

The no-enforcement stage does not attempt to redesign the local cubature rule.
Instead, it redistributes as little as possible the weights already present in
the current feasible rule.

Suppose that point $p$ carries only a small contribution in the current local
weights. Then the correction required to absorb its removal is also expected to
be small. In such cases, the redistribution remains localized and smooth. This
is precisely the regime in which the no-enforcement stage is most effective.

Conversely, when the remaining support is already too sparse, the elimination of
one point may require large compensatory redistributions. In that regime, the
unconstrained least-change update may generate:
\[
        \text{negative weights},
        \qquad
        \text{large oscillations},
        \qquad
        \text{loss of feasibility}.
\]
At that point, the algorithm transitions to the explicit positivity stage.

%-------------------------------------------------------------------------------
\subsection{Why smoothness is usually better in this stage}
%-------------------------------------------------------------------------------

The absence of active-set clamping is one of the main reasons why the
no-enforcement stage tends to preserve smoothness.

At this stage, the adaptive weights are obtained as solutions of smooth linear
least-change systems. Therefore, as long as the local operators vary smoothly
with the manifold coordinate, the resulting weight field
\[
        \qM\mapsto\wred(\qM)
\]
also tends to vary smoothly.

By contrast, once positivity is enforced, the feasible set becomes polyhedral.
As the latent coordinate changes, different constraints may become active or
inactive. These active-set transitions introduce nonsmoothness into the
solution map.

Thus, even though the no-enforcement stage allows negative weights, it often
produces more regular adaptive weight fields than the positivity-enforcing
stages.

%-------------------------------------------------------------------------------
\subsection{Acceptance and rejection of candidates}
%-------------------------------------------------------------------------------

The unconstrained redistribution associated with a candidate point $p$ is not
accepted automatically. The updated weights must satisfy several practical
criteria.

Typical checks include:
\begin{itemize}
    \item satisfaction of the local exactness systems;
    \item absence of excessive negative weights;
    \item bounded redistribution magnitude;
    \item acceptable numerical conditioning.
\end{itemize}

If the candidate satisfies the prescribed tolerances, the elimination is
committed:
\[
        \Iact
        \gets
        \Iact\setminus\{p\},
\]
and the updated adaptive weights become the new current solution:
\[
        \Wold\gets\Wnew.
\]
Otherwise, the candidate is rejected and the next point in the ordered list is
tested.

The important point is that rejection does not necessarily mean that the point
is fundamentally indispensable. It only means that the current unconstrained
redistribution stage cannot remove it while satisfying the prescribed criteria.
The same point may later become removable once positivity enforcement or
graph-regularized redistribution is activated.

%-------------------------------------------------------------------------------
\subsection{Pseudo-algorithm}
%-------------------------------------------------------------------------------

The no-enforcement redistribution stage can be summarized as follows.

\begin{algorithm}[H]
\caption{No-enforcement least-change pruning}
\label{alg:no_enforcement_pruning}
\begin{algorithmic}[1]

\Require
Current support $\Iact$;
adaptive weights $\Wold$;
local systems $\{(\Uj,\bj)\}_{j=1}^{\Nman}$.

\Ensure
Updated support and adaptive weights if a feasible elimination is found.

\State Compute aggregate weights
\[
        R_i=\sum_{j=1}^{\Nman}W_{ij}.
\]

\State Sort candidates in increasing order of $R_i$.

\For{candidate point $p$ in the ordered list}

    \State Define
    \[
        \Irem=\Iact\setminus\{p\}.
    \]

    \For{$j=1,\ldots,\Nman$}

        \State Restrict the local basis:
        \[
            \widetilde{\Uj}=\Uj(\Irem,:).
        \]

        \State Compute the least-change redistribution:
        \[
            \widetilde{\wnew}^{(j)}
            =
            \widetilde{\wold}^{(j)}
            -
            \widetilde{\Uj}
            (\widetilde{\Uj}^T\widetilde{\Uj})^{-1}
            (\widetilde{\Uj}^T\widetilde{\wold}^{(j)}-\bj).
        \]

        \State Reconstruct the full local weights with
        \[
            w_{p,\mathrm{new}}^{(j)}=0.
        \]

    \EndFor

    \State Assemble the candidate adaptive weight matrix $\Wnew$.

    \If{$\Wnew$ satisfies the acceptance criteria}

        \State Commit the elimination:
        \[
            \Iact\gets\Irem.
        \]

        \State Update
        \[
            \Wold\gets\Wnew.
        \]

        \State Exit successfully.

    \EndIf

\EndFor

\State Return failure if no candidate is accepted.

\end{algorithmic}
\end{algorithm}

%-------------------------------------------------------------------------------
\subsection{Function-level interpretation}
%-------------------------------------------------------------------------------

The no-enforcement stage is implemented in
\[
        \texttt{Loop\_MAWecmNOENF\_cl.m}.
\]
The suffix ``NOENF'' refers to the absence of explicit positivity enforcement.
The suffix ``cl'' refers historically to the least-change redistribution logic.

The function operates on all local systems simultaneously, because the support
must remain common across the manifold. However, the redistribution itself is
performed independently at each manifold node. Therefore, the stage is local in
the manifold direction and does not yet use the graph Laplacian.

Graph coupling enters only in the later graph-regularized active-set stage,
which is described in the next sections.


%===============================================================================
\section{Stage 2: local positivity enforcement by active sets}
\label{sec:local_active_set_stage}
%===============================================================================

The no-enforcement redistribution stage described previously is often capable
of removing a substantial number of integration points while preserving a
relatively smooth adaptive weight field. However, as the support becomes
progressively smaller, the unconstrained least-change redistribution may produce
negative weights or large oscillatory corrections. At that point, explicit
positivity enforcement becomes necessary.

The second redistribution stage of MAW--ECM therefore introduces nonnegativity
constraints into the local redistribution problem. This stage is implemented in
\[
        \texttt{Loop\_MAWecmNOENFe.m}.
\]
Unlike the graph-based stage described later, the present formulation still
treats each manifold node independently. The coupling between neighboring latent
states is not yet introduced.

Thus, the present stage should be viewed as:
\[
        \text{local positivity enforcement without graph coupling}.
\]

%-------------------------------------------------------------------------------
\subsection{Local constrained redistribution problem}
%-------------------------------------------------------------------------------

Suppose that a candidate point $p$ is tentatively removed from the support. As
before, define
\[
        \Irem=\Iact\setminus\{p\},
\]
and let
\[
        \widetilde{\Uj}
        =
        \Uj(\Irem,:)
\]
be the reduced local basis after elimination of the candidate point.

At manifold node $j$, the redistributed weights must satisfy:
\[
        {\widetilde{\Uj}}^T
        \widetilde{\wnew}^{(j)}
        =
        \bj.
        \label{eq:local_positive_exactness}
\]
In contrast with the no-enforcement stage, one now additionally imposes
\[
        \widetilde{\wnew}^{(j)}\geq 0.
        \label{eq:local_positive_inequality}
\]
The redistribution problem therefore becomes
\[
        \min_{\widetilde{\wnew}^{(j)}}
        \frac{1}{2}
        \|
        \widetilde{\wnew}^{(j)}
        -
        \widetilde{\wold}^{(j)}
        \|_2^2,
\]
subject to
\[
        {\widetilde{\Uj}}^T
        \widetilde{\wnew}^{(j)}
        =
        \bj,
\]
and
\[
        \widetilde{\wnew}^{(j)}\geq 0.
        \label{eq:local_positive_problem}
\]

Thus, the local adaptive weights are redistributed as little as possible while
remaining both feasible and nonnegative.

%-------------------------------------------------------------------------------
\subsection{Geometric interpretation}
%-------------------------------------------------------------------------------

Without positivity enforcement, the feasible set of the redistribution problem
is an affine subspace:
\[
        \mathcal{A}_j
        =
        \left\{
        \bm w:
        {\widetilde{\Uj}}^T\bm w=\bj
        \right\}.
\]
The no-enforcement redistribution is simply the orthogonal projection of the
previous weights onto this affine subspace.

With positivity enforcement, the feasible set becomes
\[
        \mathcal{F}_j
        =
        \left\{
        \bm w:
        {\widetilde{\Uj}}^T\bm w=\bj,
        \quad
        \bm w\geq 0
        \right\}.
\]
This set is no longer an affine space, but a convex polyhedron obtained by
intersecting the affine exactness constraints with the positive orthant.

The redistribution problem therefore becomes a constrained projection onto a
polyhedral feasible set.

%-------------------------------------------------------------------------------
\subsection{Active-set philosophy}
%-------------------------------------------------------------------------------

The implementation adopts an active-set strategy. The idea is simple: whenever
one component of the redistributed weights becomes negative, that component is
clamped to zero and treated as an additional equality constraint.

At manifold node $j$, define the local active set
\[
        \Dj
        \subset
        \{1,\ldots,|\Irem|\}.
\]
An index
\[
        i\in\Dj
\]
means that the corresponding local weight is fixed to zero:
\[
        w_i^{(j)}=0.
\]
The remaining local weights are free variables.

Thus, the constrained redistribution problem is converted into a sequence of
equality-constrained least-change problems with progressively larger active
sets.

%-------------------------------------------------------------------------------
\subsection{Reduced local variables}
%-------------------------------------------------------------------------------

Let
\[
        {\Ifree}_j
        =
        \{1,\ldots,|\Irem|\}\setminus\Dj
\]
denote the free local indices at manifold node $j$. Restricting the local basis
to the free variables gives
\[
        \widehat{\Uj}
        =
        \widetilde{\Uj}({\Ifree}_j,:).
\]
Similarly, define the free local unknown
\[
        \widehat{\wnew}^{(j)}
        =
        \widetilde{\wnew}^{(j)}({\Ifree}_j).
\]
The local constrained redistribution problem then becomes
\[
        \min_{\widehat{\wnew}^{(j)}}
        \frac{1}{2}
        \|
        \widehat{\wnew}^{(j)}
        -
        \widehat{\wold}^{(j)}
        \|_2^2,
\]
subject to
\[
        {\widehat{\Uj}}^T
        \widehat{\wnew}^{(j)}
        =
        \bj.
        \label{eq:free_active_system}
\]
This has exactly the same algebraic structure as the no-enforcement stage, but
restricted to the free variables.

%-------------------------------------------------------------------------------
\subsection{Iterative evolution of the active sets}
%-------------------------------------------------------------------------------

The active-set redistribution proceeds iteratively.

Starting from
\[
        \Dj=\emptyset,
\]
one computes the least-change redistribution over all remaining variables. If
all components are nonnegative, the candidate elimination is accepted.

Otherwise, one identifies the negative components:
\[
        \mathcal{N}_j
        =
        \left\{
        i:
        w_i^{(j)}<0
        \right\}.
\]
These indices are added to the active set:
\[
        \Dj
        \gets
        \Dj\cup\mathcal{N}_j.
\]
The redistribution problem is then solved again over the reduced free set.

This process is repeated until:
\begin{itemize}
    \item all free variables are nonnegative, or
    \item the free system becomes infeasible or singular.
\end{itemize}

Thus, the active sets evolve adaptively during the redistribution process.

%-------------------------------------------------------------------------------
\subsection{Node-dependent active sets}
%-------------------------------------------------------------------------------

A crucial aspect of the present implementation is that the active sets are local
to each manifold node.

Suppose that integration-point weight $i$ becomes negative at manifold node
$j_1$. Then
\[
        i\in{\Dj}_1,
\]
but it does not follow that
\[
        i\in{\Dj}_2
\]
for another node $j_2$.

Thus, positivity enforcement is local in latent space:
\[
        {\Dj}_1 \neq {\Dj}_2
        \quad\text{in general}.
\]
This distinction is important. The support itself remains common over the
manifold, but the local positivity constraints do not.

Consequently, even after the support has been fixed, the local constrained
redistribution problem may have different active variables at different latent
states.

This is one of the main mechanisms generating nonsmoothness in the adaptive
weight field.

%-------------------------------------------------------------------------------
\subsection{Origin of nonsmoothness}
%-------------------------------------------------------------------------------

At the no-enforcement stage, the redistribution is obtained from smooth linear
systems. Therefore, the resulting mapping
\[
        \qM\mapsto\wred(\qM)
\]
tends to vary smoothly.

Once positivity is enforced, however, the active sets may change abruptly as the
latent coordinate varies. The adaptive weight field therefore becomes the
solution map of a piecewise-defined constrained optimization problem.

Whenever the active set changes, the local affine system defining the
redistribution also changes. As a consequence:
\[
        \text{the adaptive weight field is generally only piecewise smooth}.
\]
This effect becomes increasingly pronounced as the support becomes more
compressed.

The graph-based stage introduced later is intended precisely to moderate this
behavior by coupling neighboring latent states during the redistribution.

%-------------------------------------------------------------------------------
\subsection{Acceptance and rejection}
%-------------------------------------------------------------------------------

A candidate point $p$ is accepted only if the positivity-enforcing active-set
procedure succeeds at every manifold node.

More precisely, for all
\[
        j=1,\ldots,\Nman,
\]
the final redistributed weights must satisfy:
\[
        {\widetilde{\Uj}}^T
        \widetilde{\wnew}^{(j)}
        =
        \bj,
\]
and
\[
        \widetilde{\wnew}^{(j)}\geq 0,
\]
up to the prescribed tolerances.

If the constrained redistribution fails at even one manifold node, the
candidate elimination is rejected.

Thus, feasibility is global with respect to the sampled manifold, even though
the redistributions themselves are computed locally.

%-------------------------------------------------------------------------------
\subsection{Pseudo-algorithm}
%-------------------------------------------------------------------------------

The local active-set redistribution can be summarized as follows.

\begin{algorithm}[H]
\caption{Local positivity-enforcing redistribution}
\label{alg:local_active_set}
\begin{algorithmic}[1]

\Require
Candidate point $p$;
current adaptive weights $\Wold$;
local systems $\{(\Uj,\bj)\}_{j=1}^{\Nman}$.

\Ensure
Updated adaptive weights if the candidate is feasible.

\State Define
\[
        \Irem=\Iact\setminus\{p\}.
\]

\For{$j=1,\ldots,\Nman$}

    \State Initialize
    \[
        \Dj=\emptyset.
    \]

    \While{active-set iteration has not converged}

        \State Define the free local indices
        \[
            {\Ifree}_j
            =
            \{1,\ldots,|\Irem|\}\setminus\Dj.
        \]

        \State Restrict the basis to the free variables
        \[
            \widehat{\Uj}
            =
            \widetilde{\Uj}({\Ifree}_j,:).
        \]

        \State Solve the free least-change redistribution problem.

        \If{all free variables are nonnegative}

            \State Accept the local redistribution.
            \State Exit the local active-set loop.

        \Else

            \State Detect negative free variables.

            \State Add them to the active set:
            \[
                \Dj
                \gets
                \Dj\cup\mathcal{N}_j.
            \]

        \EndIf

    \EndWhile

    \If{the local redistribution failed}

        \State Reject the candidate point $p$ globally.
        \State Exit with failure.

    \EndIf

\EndFor

\State Accept the candidate point $p$ globally.

\end{algorithmic}
\end{algorithm}

%-------------------------------------------------------------------------------
\subsection{Interpretation}
%-------------------------------------------------------------------------------

The local active-set stage can be viewed as a compromise between two competing
requirements.

On one hand, the adaptive weights should remain close to the previous feasible
solution:
\[
        \Wnew \approx \Wold.
\]
On the other hand, positivity must be enforced:
\[
        \Wnew\geq 0.
\]
As the support becomes smaller, these two requirements become increasingly
difficult to reconcile simultaneously.

The resulting constrained redistribution is therefore intrinsically nonlinear
and nonsmooth. The graph-based stage described next attempts to compensate for
this by coupling neighboring manifold states and discouraging abrupt local
redistributions.

%-------------------------------------------------------------------------------
\subsection{Function-level interpretation}
%-------------------------------------------------------------------------------

The local positivity-enforcing stage is implemented in
\[
        \texttt{Loop\_MAWecmNOENFe.m}.
\]
The suffix ``NOENFe'' historically refers to the no-overlap local active-set
version with explicit enforcement.

The key structural property of this routine is:
\[
        \text{local active sets, but common support}.
\]
The support elimination is global, but the positivity constraints are enforced
independently at each manifold node.

This distinction will become important in the graph-based stage, where the
redistribution is instead performed through a coupled optimization problem over
all manifold states simultaneously.


%===============================================================================
\section{Stage 3: graph-regularized positivity enforcement}
\label{sec:graph_regularized_stage}
%===============================================================================

The local active-set stage enforces positivity, but it treats each manifold
sample independently. This independence is convenient algebraically, yet it may
produce weight fields with sharp variations across neighboring latent states.
The graph-regularized stage is introduced to reduce this effect.

This stage is implemented through
\[
        \texttt{GraphBased\_MAW\_ECM\_2v.m},
\]
which calls the core redistribution routine
\[
        \texttt{prune\_step\_optionB.m}.
\]
The essential change with respect to the local active-set stage is that the
weights at all manifold samples are updated simultaneously. The graph Laplacian
$\Kgraph$ couples neighboring latent nodes and penalizes rough redistributions.

%-------------------------------------------------------------------------------
\subsection{Global unknown and constraints}
%-------------------------------------------------------------------------------

Let the current adaptive weights be
\[
        \Wold \in \R^{r\times\Nman},
\]
where $r$ denotes the number of currently retained candidate points at this
stage of the pruning loop. A tentative elimination of point $p$ seeks updated
weights
\[
        \Wnew \in \R^{r\times\Nman}.
\]
The pruning condition is imposed globally:
\[
        W^{\mathrm{new}}_{pj}=0,
        \qquad
        j=1,\ldots,\Nman.
        \label{eq:graph_pruned_row}
\]
For each manifold node $j$, the local exactness condition remains
\[
        A_j {\wnew}^{(j)} = \bj,
        \label{eq:graph_local_exactness}
\]
where, in the notation of the implementation,
\[
        A_j = {\Uj}^T
\]
restricted to the rows currently active in the pruning trial.

The positivity condition is
\[
        {\wnew}^{(j)}\geq 0,
        \qquad
        j=1,\ldots,\Nman.
        \label{eq:graph_positivity}
\]

%-------------------------------------------------------------------------------
\subsection{Incremental graph-regularized objective}
%-------------------------------------------------------------------------------

The graph-based update minimizes a least-change functional augmented with a
smoothness penalty over the latent graph. In incremental form, the objective is
\[
        \frac{1}{2}
        \|\Wnew-\Wold\|_F^2
        +
        \frac{\alphaG}{2}
        \trace\left[
        (\Wnew-\Wold)
        \Kgraph
        (\Wnew-\Wold)^T
        \right].
        \label{eq:graph_regularized_objective}
\]
The first term keeps the new adaptive weights close to the previous feasible
weights. The second term penalizes rough variations of the update over the
sampled latent manifold.

The use of the increment
\[
        \DeltaW=\Wnew-\Wold
\]
is important. It means that the graph regularization acts on the redistribution
caused by the pruning step, not necessarily on the absolute value of the
weights. Thus, a weight field that is already nonsmooth is not directly forced
to become smooth in one step; rather, the new pruning-induced change is
discouraged from introducing additional sharp variations.

If the absolute version is used instead, the graph term is applied directly to
$\Wnew$. The incremental version is usually less invasive and better aligned
with the idea of progressive pruning.

%-------------------------------------------------------------------------------
\subsection{Stacked quadratic form}
%-------------------------------------------------------------------------------

The implementation solves the graph-regularized redistribution in stacked-vector
form. Define
\[
        \zvec=\vecop(\Wnew),
        \qquad
        \zold=\vecop(\Wold),
\]
with columnwise stacking. Then
\[
        \|\Wnew-\Wold\|_F^2
        =
        \|\zvec-\zold\|_2^2,
\]
and
\[
        \trace\left[
        (\Wnew-\Wold)\Kgraph(\Wnew-\Wold)^T
        \right]
        =
        (\zvec-\zold)^T
        (\Kgraph\otimes\bm I_r)
        (\zvec-\zold).
\]
Therefore, the objective can be written as
\[
        \frac{1}{2}
        (\zvec-\zold)^T
        \Hgraph
        (\zvec-\zold),
        \label{eq:graph_stacked_incremental}
\]
where
\[
        \Hgraph
        =
        \bm I
        +
        \alphaG
        (\Kgraph\otimes\bm I_r).
        \label{eq:graph_hessian}
\]
Expanding the quadratic form and dropping constants gives
\[
        \frac{1}{2}\zvec^T\Hgraph\zvec
        -
        \bm g^T\zvec,
        \label{eq:graph_quadratic_expanded}
\]
with
\[
        \bm g = \Hgraph\zold
\]
in the incremental formulation.

The matrix $\Hgraph$ is symmetric positive definite for $\alphaG\geq 0$, since
$\Kgraph$ is positive semidefinite and the identity term is present. This is
important for the robustness of the reduced nullspace solve.

%-------------------------------------------------------------------------------
\subsection{Node-dependent active sets}
%-------------------------------------------------------------------------------

The graph-based implementation enforces nonnegativity through node-dependent
active sets. For each manifold node $j$, define
\[
        \Dj \subset \{1,\ldots,r\}.
\]
The interpretation is:
\[
        i\in\Dj
        \quad\Longleftrightarrow\quad
        W^{\mathrm{new}}_{ij}=0.
\]
The pruning candidate is clamped at all nodes:
\[
        p\in\Dj,
        \qquad
        j=1,\ldots,\Nman.
\]
Additional indices are added to $\Dj$ only at the nodes where negative weights
are detected.

This is the key distinction of the option implemented in
\texttt{prune\_step\_optionB.m}: a weight can be clamped at one manifold node
without being clamped at all other nodes. Hence, the active set is local in the
latent direction, even though the update is graph-coupled.

%-------------------------------------------------------------------------------
\subsection{Nullspace parametrization}
%-------------------------------------------------------------------------------

For a fixed collection of active sets
\[
        \{\Dj\}_{j=1}^{\Nman},
\]
the problem contains only equality constraints: local exactness constraints and
clamping constraints. The implementation solves this equality-constrained
quadratic problem by nullspace parametrization.

Let
\[
        \zvec = \zp + \Nmat\bm y,
        \label{eq:nullspace_parametrization}
\]
where:
\begin{itemize}
    \item $\zp$ is one particular vector satisfying all equality constraints;
    \item $\Nmat$ is a basis for the nullspace of the equality-constraint
    matrix;
    \item $\bm y$ contains the reduced unknowns.
\end{itemize}

Substitution into the expanded quadratic objective gives the reduced system
\[
        \Hred \bm y = \fred,
        \label{eq:reduced_nullspace_system}
\]
where
\[
        \Hred = \Nmat^T\Hgraph\Nmat,
\]
and
\[
        \fred =
        \Nmat^T
        \left(
        \bm g-\Hgraph\zp
        \right).
\]
Since $\Hgraph$ is positive definite and $\Nmat$ has full column rank, the
matrix $\Hred$ is also positive definite. The reduced system can therefore be
solved directly.

The full updated weight vector is then reconstructed from
\[
        \zvec = \zp+\Nmat\bm y,
\]
and reshaped into
\[
        \Wnew\in\R^{r\times\Nman}.
\]

%-------------------------------------------------------------------------------
\subsection{Active-set iterations}
%-------------------------------------------------------------------------------

After solving the equality-constrained graph-regularized problem for the
current active sets, the algorithm checks nonnegativity. For each node $j$, the
free variables are inspected. If
\[
        W^{\mathrm{new}}_{ij}<-\varepsilon_{\mathrm{neg}},
\]
for some free index $i$, that index is added to the local active set:
\[
        \Dj \gets \Dj\cup\{i\}.
\]
The graph-coupled problem is then solved again with the enlarged active sets.

The active-set iteration stops when no new negative weights are detected. If the
maximum number of active-set iterations is reached, or if the constrained system
becomes infeasible, the pruning candidate is rejected.

Thus, one graph-based pruning trial consists of:
\[
        \text{fix candidate }p
        \;\longrightarrow\;
        \text{solve graph-coupled equality problem}
        \;\longrightarrow\;
        \text{update active sets}
        \;\longrightarrow\;
        \text{repeat until feasible or rejected}.
\]

%-------------------------------------------------------------------------------
\subsection{Why graph coupling helps}
%-------------------------------------------------------------------------------

In the local active-set stage, each manifold node redistributes its weights
independently. If a weight is clamped at one node but not at its neighbors, the
resulting field may develop a localized jump. The graph-based stage reduces this
effect by penalizing large differences of the update across neighboring latent
nodes.

For an edge $(a,b)\in\Egraph$, the update jump associated with integration point
$i$ is
\[
        \Delta W_{ia}-\Delta W_{ib}.
\]
The graph energy penalizes such jumps in an averaged quadratic sense. Therefore,
among the feasible redistributions, the graph-based update favors those with
more coherent changes across the latent manifold.

This does not eliminate active-set nonsmoothness. If the constraints and
positivity conditions force a local clamp, the graph penalty cannot remove it.
It can only choose, among the feasible solutions, the one that is least rough
according to the chosen graph metric.

%-------------------------------------------------------------------------------
\subsection{Pseudo-algorithm}
%-------------------------------------------------------------------------------

\begin{algorithm}[H]
\caption{Graph-regularized active-set redistribution for one pruning candidate}
\label{alg:graph_regularized_pruning}
\begin{algorithmic}[1]

\Require
Current weights $\Wold$;
local systems $\{(A_j,\bj)\}_{j=1}^{\Nman}$;
graph Laplacian $\Kgraph$;
smoothing parameter $\alphaG$;
candidate point $p$.

\Ensure
Updated weights $\Wnew$ if the pruning candidate is feasible.

\State Initialize active sets:
\[
        \Dj\gets\{p\},
        \qquad
        j=1,\ldots,\Nman.
\]

\State Build
\[
        \Hgraph
        =
        \bm I+\alphaG(\Kgraph\otimes\bm I_r).
\]

\If{incremental smoothing is used}
    \State Set
    \[
        \bm g=\Hgraph\zold.
    \]
\Else
    \State Set
    \[
        \bm g=\zold.
    \]
\EndIf

\For{active-set iteration $k=1,\ldots,k_{\max}$}

    \State Build the global equality constraints from:
    \[
        A_j{\wnew}^{(j)}=\bj,
        \qquad
        W^{\mathrm{new}}_{ij}=0
        \quad
        \text{for } i\in\Dj.
    \]

    \State Compute a particular solution $\zp$ and nullspace basis $\Nmat$.

    \State Solve
    \[
        \left(\Nmat^T\Hgraph\Nmat\right)\bm y
        =
        \Nmat^T(\bm g-\Hgraph\zp).
    \]

    \State Reconstruct
    \[
        \zvec=\zp+\Nmat\bm y,
        \qquad
        \Wnew=\operatorname{reshape}(\zvec).
    \]

    \State Detect negative free weights.

    \If{no new negative weights are detected}
        \State Accept the candidate and return success.
    \Else
        \State Add the negative entries to the corresponding node-dependent
        active sets $\Dj$.
    \EndIf

\EndFor

\State Reject the candidate if convergence is not achieved.

\end{algorithmic}
\end{algorithm}

%-------------------------------------------------------------------------------
\subsection{Function-level interpretation}
%-------------------------------------------------------------------------------

At the orchestration level, graph-based pruning is performed by
\[
        \texttt{GraphBased\_MAW\_ECM\_2v.m}.
\]
This routine tests pruning candidates, calls the graph-regularized update, and
selects one candidate to commit.

The actual graph-regularized active-set solve for one candidate is implemented
in
\[
        \texttt{prune\_step\_optionB.m}.
\]
This routine receives the current weights, the local exactness systems, the
graph Laplacian, the smoothing parameter, and the candidate index. It returns
the updated weights and diagnostic information indicating whether the candidate
was feasible.

The name ``Option B'' refers to the node-dependent active-set strategy:
negative weights are clamped only at the manifold nodes where they occur, rather
than globally across all nodes.


%===============================================================================
\section{Regression of the adaptive weight fields}
\label{sec:weight_regression}
%===============================================================================

After the pruning stage has terminated, MAW--ECM provides:
\[
        \Zred
        \subset
        \Zfull,
\]
together with the sampled adaptive weights
\[
        \Wadapt_{\mathrm{red}}
        \in
        \R^{\Nred\times\Nman}.
\]
The columns of $\Wadapt_{\mathrm{red}}$ correspond to the sampled latent
coordinates
\[
        \{\qMj\}_{j=1}^{\Nman}.
\]
At this stage, the adaptive weights are known only at the sampled manifold
states used during the offline stage. An additional approximation step is
therefore required in order to evaluate the adaptive cubature rule at new
online latent coordinates.

The purpose of the present stage is to construct the mapping
\[
        \qM
        \longmapsto
        \wred(\qM).
\]
This stage is implemented in
\[
        \texttt{RegressionViaRBF\_master\_weights.m}.
\]

%-------------------------------------------------------------------------------
\subsection{Weight fields over the latent manifold}
%-------------------------------------------------------------------------------

Each retained integration point defines one scalar field over the latent
manifold. Let
\[
        \bm w_i
        =
        \begin{bmatrix}
        W_{i1} &
        W_{i2} &
        \cdots &
        W_{i\Nman}
        \end{bmatrix}^T
        \in\R^{\Nman}
\]
collect the adaptive weights associated with retained integration point $i$
over all sampled manifold states.

Thus, the adaptive cubature rule is not represented by a single vector of
weights, but by a family of scalar fields:
\[
        \{
        w_i(\qM)
        \}_{i=1}^{\Nred}.
\]
The regression stage approximates these scalar fields from the sampled data
\[
        \left\{
        \qMj,
        \;
        W_{ij}
        \right\}_{j=1}^{\Nman}.
\]

%-------------------------------------------------------------------------------
\subsection{Why regression is necessary}
%-------------------------------------------------------------------------------

The pruning stage operates only on the sampled manifold states available during
offline training. However, during online evaluation, the latent coordinates
will generally not coincide with the sampled nodes:
\[
        {\qM}_{\mathrm{online}}
        \notin
        \{
        \qMj
        \}_{j=1}^{\Nman}.
\]
Therefore, the adaptive weights must be interpolated or approximated.

The online adaptive cubature rule takes the form
\[
        I({\qM}_{\mathrm{online}})
        \approx
        \sum_{i\in\Zred}
        f_i({\qM}_{\mathrm{online}})
        \,
        w_i({\qM}_{\mathrm{online}}).
\]
Hence, online evaluation requires the ability to compute
\[
        \wred({\qM}_{\mathrm{online}}).
\]

The regression stage should therefore be interpreted as the final step that
converts the sampled adaptive weights into a continuous surrogate model over
the latent manifold.

%-------------------------------------------------------------------------------
\subsection{Choice of radial basis functions}
%-------------------------------------------------------------------------------

The present implementation uses radial basis functions (RBFs). This choice is
motivated by several considerations.

First, the latent dimension is very small:
\[
        \Nlat=2.
\]
Second, the sampled latent manifold is structured and relatively regular.
Third, the adaptive weights, although not perfectly smooth, usually vary in a
coherent manner over the manifold after graph regularization.

Under these conditions, RBF interpolation provides:
\begin{itemize}
    \item meshfree approximation,
    \item smooth evaluation,
    \item simple implementation,
    \item differentiability of the resulting surrogate,
    \item flexibility with respect to irregular sampling.
\end{itemize}

The current implementation uses anisotropic Gaussian kernels.

%-------------------------------------------------------------------------------
\subsection{RBF approximation of one weight field}
%-------------------------------------------------------------------------------

Consider one retained integration point $i$. The corresponding scalar weight
field is approximated as
\[
        w_i(\qM)
        \approx
        \sum_{k=1}^{n_{\mathrm{RBF}}}
        \gamma_{ik}
        \,
        \phi_k(\qM)
        +
        p_i(\qM),
        \label{eq:rbf_weight_approximation}
\]
where:
\begin{itemize}
    \item $\phi_k$ are radial basis functions centered at selected latent
    samples,
    \item $\gamma_{ik}$ are the regression coefficients,
    \item $p_i(\qM)$ is a low-order polynomial correction.
\end{itemize}

For Gaussian kernels,
\[
        \phi_k(\qM)
        =
        \exp
        \left(
        -\frac{1}{2}
        (\qM-\bm c_k)^T
        \bm D_k^{-2}
        (\qM-\bm c_k)
        \right),
        \label{eq:gaussian_rbf}
\]
where:
\begin{itemize}
    \item $\bm c_k$ is the center of the $k$-th RBF,
    \item $\bm D_k$ defines anisotropic characteristic lengths.
\end{itemize}

The polynomial correction may include:
\[
        p_i(\qM)
        =
        a_i,
        \qquad
        \text{or}
        \qquad
        p_i(\qM)
        =
        a_i+\bm b_i^T\qM,
\]
depending on the selected approximation order.

%-------------------------------------------------------------------------------
\subsection{Structured subsampling of RBF centers}
%-------------------------------------------------------------------------------

Using all sampled manifold states as RBF centers is often unnecessary and may
lead to large dense interpolation systems. The implementation therefore performs
structured subsampling of the latent grid.

Let
\[
        \{
        \bm c_k
        \}_{k=1}^{n_{\mathrm{RBF}}}
\]
denote the retained centers. Typically,
\[
        n_{\mathrm{RBF}}
        \ll
        \Nman.
\]
The centers are selected so as to preserve coverage of the latent domain while
reducing redundancy among nearby samples.

This choice is especially useful because neighboring latent states usually carry
very similar adaptive weights after graph regularization.

%-------------------------------------------------------------------------------
\subsection{Matrix form of the regression problem}
%-------------------------------------------------------------------------------

For one integration point $i$, define the sampled values
\[
        \bm w_i
        =
        \begin{bmatrix}
        W_{i1}\\
        \vdots\\
        W_{i\Nman}
        \end{bmatrix}.
\]
Define also the interpolation matrix
\[
        \bm\Phi
        \in
        \R^{\Nman\times n_{\mathrm{RBF}}},
\]
with entries
\[
        \Phi_{jk}
        =
        \phi_k(\qMj).
\]
If polynomial correction is included, let
\[
        \bm P
\]
denote the polynomial basis evaluated at the sampled latent coordinates.

The regression problem can then be written schematically as
\[
        \begin{bmatrix}
        \bm\Phi & \bm P
        \end{bmatrix}
        \begin{bmatrix}
        \bm\gamma_i\\
        \bm\eta_i
        \end{bmatrix}
        \approx
        \bm w_i.
        \label{eq:rbf_linear_system}
\]
Depending on the selected options, the implementation may use:
\begin{itemize}
    \item interpolation,
    \item least-squares approximation,
    \item regularized least squares.
\end{itemize}

The regression is performed independently for each retained integration point.

%-------------------------------------------------------------------------------
\subsection{Differentiability}
%-------------------------------------------------------------------------------

One important consequence of the RBF representation is that the adaptive
weights become differentiable functions of the latent coordinates:
\[
        \wred(\qM)
        \in
        C^\infty
\]
for Gaussian kernels.

Thus, one may compute:
\[
        \frac{\partial \wred}{\partial \qM},
        \qquad
        \frac{\partial^2 \wred}{\partial \qM^2}.
\]
This property may become useful in future extensions involving:
\begin{itemize}
    \item tangent-consistent hyperreduction,
    \item latent-space continuation,
    \item adaptive online integration,
    \item reduced Newton schemes,
    \item sensitivity analysis.
\end{itemize}

It should be emphasized, however, that the smoothness of the RBF surrogate does
not imply that the underlying sampled weights are themselves smooth. If the
offline adaptive weights contain strong active-set transitions, the RBF
approximation simply produces a smooth interpolation of a piecewise-smooth
dataset.

%-------------------------------------------------------------------------------
\subsection{Influence of graph regularization}
%-------------------------------------------------------------------------------

The graph-regularized pruning stage plays an important role in the quality of
the regression problem.

Without graph regularization, neighboring manifold samples may carry strongly
different adaptive weights due to local active-set changes. In that case, the
mapping
\[
        \qM\mapsto\wred(\qM)
\]
becomes more difficult to regress accurately.

The graph penalty mitigates this effect by discouraging abrupt redistribution
patterns across neighboring latent states. Consequently:
\begin{itemize}
    \item the sampled adaptive weights become more coherent,
    \item the regression problem becomes better conditioned,
    \item fewer RBF centers are generally required,
    \item online interpolation becomes more stable.
\end{itemize}

Thus, graph regularization should be viewed not only as a stabilization device
for pruning, but also as a preprocessing step for the final surrogate model of
the adaptive weights.

%-------------------------------------------------------------------------------
\subsection{Online evaluation}
%-------------------------------------------------------------------------------

Once the regression coefficients have been computed, online evaluation proceeds
as follows.

Given a latent coordinate
\[
        {\qM}_{\mathrm{online}},
\]
the RBF surrogate evaluates
\[
        \wred({\qM}_{\mathrm{online}})
        \in
        \R^{\Nred}.
\]
The online adaptive cubature approximation is then
\[
        I({\qM}_{\mathrm{online}})
        \approx
        \sum_{i\in\Zred}
        f_i({\qM}_{\mathrm{online}})
        \,
        w_i({\qM}_{\mathrm{online}}).
\]
Thus, online evaluation requires:
\begin{itemize}
    \item evaluation of the reduced set of integrands,
    \item evaluation of the adaptive weights through the RBF surrogate.
\end{itemize}

Importantly, the online stage does not solve a new ECM problem. All adaptive
behavior has already been encoded offline in the learned mapping
\[
        \qM\mapsto\wred(\qM).
\]

%-------------------------------------------------------------------------------
\subsection{Pseudo-algorithm}
%-------------------------------------------------------------------------------

\begin{algorithm}[H]
\caption{Regression of the adaptive weight fields}
\label{alg:weight_regression}
\begin{algorithmic}[1]

\Require
Sampled latent coordinates
\[
        \{\qMj\}_{j=1}^{\Nman};
\]
sampled adaptive weights
\[
        \Wadapt_{\mathrm{red}};
\]
RBF parameters and regression options.

\Ensure
Surrogate model
\[
        \qM\mapsto\wred(\qM).
\]

\State Select the RBF centers
\[
        \{
        \bm c_k
        \}_{k=1}^{n_{\mathrm{RBF}}}.
\]

\State Build the interpolation matrix
\[
        \Phi_{jk}=\phi_k(\qMj).
\]

\For{each retained integration point $i=1,\ldots,\Nred$}

    \State Extract the sampled weight field
    \[
        \bm w_i=
        \begin{bmatrix}
        W_{i1}\\
        \vdots\\
        W_{i\Nman}
        \end{bmatrix}.
    \]

    \State Solve the RBF regression problem.

    \State Store the coefficients:
    \[
        \bm\gamma_i,
        \qquad
        \bm\eta_i.
    \]

\EndFor

\State Return the RBF surrogate model.

\end{algorithmic}
\end{algorithm}

%-------------------------------------------------------------------------------
\subsection{Function-level interpretation}
%-------------------------------------------------------------------------------

The regression stage is implemented in
\[
        \texttt{RegressionViaRBF\_master\_weights.m}.
\]
This routine receives:
\begin{itemize}
    \item the sampled latent coordinates,
    \item the adaptive weights,
    \item the RBF interpolation parameters.
\end{itemize}

The resulting regression structure is then evaluated online through routines of
the same family as:
\[
        \texttt{EvalSlaveRBF\_Online.m},
\]
although specialized to the adaptive weight fields.

The present implementation is specifically tuned for:
\[
        \Nlat=2,
\]
structured latent grids, and visualization of the adaptive weight surfaces.
However, the regression formulation itself is not restricted to the bivariate
case. The same approximation strategy extends naturally to higher-dimensional
latent manifolds, provided that:
\begin{itemize}
    \item the latent coordinates remain sufficiently low-dimensional,
    \item the sampled manifold is adequately covered,
    \item the adaptive weights vary coherently over the latent space.
\end{itemize}


%===============================================================================
\section{A physical interpretation of the pruning process}
\label{sec:physical_interpretation}
%===============================================================================

Before continuing with the mathematical details of the redistribution and
active-set procedures, it is useful to pause and develop a more physical and
perhaps more appealing interpretation of the MAW--ECM mechanism.

The goal of the present section is not mathematical rigor. Rather, the idea is
to help the reader build an intuitive mental picture of what the algorithm is
actually doing. In practice, many aspects of the method become much easier to
understand once one stops thinking in terms of matrices and optimization
variables and instead imagines a physical structure evolving under multiple
loading conditions.

The analogy that follows should therefore be understood mainly as a conceptual
device intended to facilitate understanding of the method without entering
immediately into all the subtleties of constrained optimization.

%-------------------------------------------------------------------------------
\subsection{A suspended bridge under many wind conditions}
%-------------------------------------------------------------------------------

Imagine a very large suspended bridge composed of a dense network of tensile
cables. Suppose now that the bridge must remain in equilibrium not under one
single loading condition, but under an entire family of nearby conditions:
slightly different wind directions, slightly different traffic distributions,
slightly different deformation states, and so on.

This is the role played by the manifold samples
\[
        \qMj.
\]
Each manifold point may be interpreted as one loading configuration of the same
bridge.

Now imagine that, initially, the bridge has been designed in a very conservative
way. There are many more cables than are truly necessary. The structure is
highly redundant.

However, not all cables play the same role.

Some cables are absolutely crucial. They carry large portions of the load and
participate in balancing many equilibrium conditions simultaneously. Others are
less important and contribute only marginally.

This is precisely what one observes in the initial ECM rule.

Empirically, the ECM weights are highly nonuniform:
some weights become much larger than others. This is not accidental. It is a
direct consequence of the greedy nature of ECM. The algorithm progressively
selects those integration points that are most effective for reproducing the
reduced equilibrium conditions.

Under the structural analogy, the ECM algorithm has already identified the
principal load paths of the bridge.

The resulting structure is therefore not homogeneous. Some cables carry large
tensions, while others play only secondary corrective roles.

%-------------------------------------------------------------------------------
\subsection{The frozen-force philosophy of standard ECM}
%-------------------------------------------------------------------------------

At this point, one may imagine the philosophy of a standard fixed-weight ECM
rule in the following way.

Suppose the engineer designing the bridge imposes the following artificial
restriction:
\begin{quote}
every cable must always carry exactly the same tension pattern for all loading
configurations.
\end{quote}

Immediately, the bridge becomes much harder to optimize.

Why?

Because different loading conditions naturally require different stress
redistributions. A cable that is extremely important under one wind direction
may become almost irrelevant under another.

Nevertheless, the fixed-weight ECM philosophy forces the same force
distribution to remain valid everywhere on the manifold.

As a consequence, the bridge requires many more cables than are locally
necessary under individual loading states.

This is precisely the main limitation of a fixed manifold ECM rule.

%-------------------------------------------------------------------------------
\subsection{The central idea of MAW--ECM}
%-------------------------------------------------------------------------------

The central idea of MAW--ECM is almost deceptively simple.

Instead of forcing the cable tensions to remain frozen over all loading
conditions, one allows the internal force distribution to adapt to the loading
configuration.

The bridge keeps the same surviving cables:
\[
        \text{same support},
\]
but the tensions carried by those cables may vary from one loading state to
another:
\[
        \text{adaptive weights}.
\]

Thus, the adaptive weights
\[
        \wadapt^{(j)}
\]
may be interpreted as the tensions carried by the surviving cables under the
loading configuration associated with manifold point $\qMj$.

This immediately introduces additional flexibility into the system.

A cable that is critically important for one loading condition may carry very
little tension under another. The structure can now reorganize its internal
forces locally over the manifold.

%-------------------------------------------------------------------------------
\subsection{Pruning as cutting cables}
%-------------------------------------------------------------------------------

Now imagine beginning a process of progressive simplification.

One cable is selected and cut.

Immediately, the external loading conditions remain exactly the same. The wind
does not disappear merely because one cable has been removed. Therefore, the
remaining cables must reorganize their tensions so as to preserve equilibrium.

At the beginning, this is easy.

The bridge is highly redundant. There are many alternative equilibrium paths.
When one cable disappears, neighboring cables simply increase their tensions
slightly and the structure remains perfectly stable.

This is the meaning of the null space in the MAW--ECM formulation.

The null-space directions represent admissible internal force redistributions
that preserve the same external equilibrium resultants. The structure can
``wiggle'' its internal stresses while keeping the same global balance.

However, every removed cable reduces the redundancy of the structure.

The remaining bridge becomes progressively more stressed:
\begin{itemize}
    \item fewer cables remain;
    \item each cable carries more responsibility;
    \item the redistribution freedom decreases;
    \item the null space shrinks.
\end{itemize}

Eventually, the bridge approaches a critical regime in which the surviving
cables are almost uniquely determined by equilibrium.

%-------------------------------------------------------------------------------
\subsection{Slacking cables and active sets}
%-------------------------------------------------------------------------------

Now comes the most important part of the analogy.

Suppose that, after removing another cable, the redistribution problem attempts
to assign a negative tension to one of the surviving members.

Mechanically, this means:
\begin{quote}
the structure is asking a cable to push.
\end{quote}

But a cable cannot push.

A cable can only pull.

This is precisely the meaning of a negative weight in MAW--ECM.

The positivity constraint
\[
        W_{ij}\ge0
\]
therefore expresses a tension-only admissibility condition.

The active-set mechanism then acquires a completely natural interpretation.

Whenever a cable attempts to enter compression, it simply goes slack:
\[
        W_{ij}=0.
\]

Importantly, this does not necessarily remove the cable globally from the
bridge. The same cable may:
\begin{itemize}
    \item carry tension under some loading configurations,
    \item become slack under others,
    \item and reactivate elsewhere on the manifold.
\end{itemize}

This distinction is fundamental.

Pruning corresponds to permanently cutting a cable from the bridge.

The active set, by contrast, determines which surviving cables are actually
load-bearing under each local loading configuration.

%-------------------------------------------------------------------------------
\subsection{Regularization as coherent structural behavior}
%-------------------------------------------------------------------------------

Finally, the role of regularization also becomes very intuitive.

Neighboring manifold points correspond to nearby loading configurations. One
therefore expects nearby loading conditions to produce nearby internal stress
patterns.

Without regularization, however, tiny changes in the loading state may produce
wildly different redistributions:
\begin{itemize}
    \item one cable active here,
    \item slack nearby,
    \item active again elsewhere.
\end{itemize}

The structure would behave erratically.

The regularization discourages this behavior. It favors coherent stress
evolution throughout the manifold and promotes smoother transitions between
neighboring loading configurations.

%-------------------------------------------------------------------------------
\subsection{Final picture}
%-------------------------------------------------------------------------------

Under this interpretation, MAW--ECM may be viewed as a progressive search for
the smallest admissible tensile skeleton capable of supporting an entire family
of neighboring equilibrium states.

The initial ECM rule provides an already compressed bridge obtained through a
greedy equilibrium process. The pruning stage progressively cuts cables from
the structure. The remaining cables continuously reorganize their tensions so
as to preserve equilibrium. The null space measures the remaining redundancy of
the bridge. The active set describes locally slack cables. The regularization
promotes coherent structural behavior over the manifold.

Although purely conceptual, this picture captures surprisingly well the true
mechanism underlying the MAW--ECM redistribution process.


\end{document}
